\documentclass[11pt,a4paper]{article}
\usepackage[T1]{fontenc}
\usepackage[utf8]{inputenc}
\usepackage{lmodern,microtype}
\usepackage[margin=27mm]{geometry}
\usepackage{amsmath,amssymb,amsthm,mathtools}
\usepackage{booktabs,longtable,array,enumitem}
\usepackage{needspace}
\usepackage{tikz-cd}
\usepackage{xcolor}
\usepackage[colorlinks=true,linkcolor=blue!55!black,citecolor=blue!55!black,urlcolor=blue!55!black]{hyperref}
\usepackage{bookmark}
\hypersetup{pdftitle={A first introduction to Deligne--Mumford stacks},pdfauthor={Notes prepared for the AlgebraicCutting project}}
\setlength{\parskip}{4pt}
\setlength{\parindent}{0pt}
\setlength{\emergencystretch}{2em}
\setlist{itemsep=3pt,topsep=5pt}
\numberwithin{equation}{section}
\newtheorem{proposition}{Proposition}[section]
\newtheorem{theorem}[proposition]{Theorem}
\newtheorem{definition}[proposition]{Definition}
\newtheorem{lemma}[proposition]{Lemma}
\theoremstyle{definition}
\newtheorem{example}[proposition]{Example}
\newtheorem{exercise}{Exercise}
\newtheorem{remark}[proposition]{Remark}
\newcommand{\A}{\mathbb A}
\newcommand{\PP}{\mathbb P}
\newcommand{\Gm}{\mathbb G_m}
\newcommand{\OO}{\mathcal O}
\newcommand{\cX}{\mathcal X}
\newcommand{\cY}{\mathcal Y}
\newcommand{\cZ}{\mathcal Z}
\newcommand{\cD}{\mathcal D}
\newcommand{\cL}{\mathcal L}
\newcommand{\cM}{\mathcal M}
\newcommand{\cP}{\mathcal P}
\newcommand{\Spec}{\operatorname{Spec}}
\newcommand{\Proj}{\operatorname{Proj}}
\newcommand{\Hom}{\operatorname{Hom}}
\newcommand{\Aut}{\operatorname{Aut}}
\newcommand{\Isom}{\operatorname{Isom}}
\newcommand{\Stab}{\operatorname{Stab}}
\newcommand{\Pic}{\operatorname{Pic}}
\newcommand{\sst}{\mathrm{ss}}
\newcommand{\id}{\mathrm{id}}
\newcommand{\one}{\mathbf 1}
\newcommand{\normal}{\mathcal N}
\newcommand{\charline}[1]{k_{#1}}
\newcommand{\reading}[1]{\par\smallskip\textit{Reading checkpoint.} #1\par}
\newcommand{\source}[1]{\par\smallskip{\small\textit{Reference.} #1}\par}
\title{A first introduction to Deligne--Mumford stacks\\[5pt]
\large Quotients, families, line bundles, and algebraic cutting}
\author{Notes prepared for the AlgebraicCutting project}
\date{13 November 2026}

\begin{document}
\maketitle
\begin{abstract}
These notes start from schemes, gluing, and line bundles, at the level of
Hartshorne, Chapters I--II. They explain why quotient spaces can lose geometric
information, how groupoids and descent retain it, and what the definition of a
Deligne--Mumford stack says. A small collection of examples is developed repeatedly:
the classifying stack $B\mu_n$, the quotient $[\A^1/\mu_n]$, a quadratic cone,
and the weighted projective line $\PP(n,1)$. The last example is also an explicit
algebraic cut. It makes the distinction between a stack line bundle, its
pushforward to the coarse space, and a line bundle descended to that space visible
in elementary calculations. Definitions and worked arguments are separated from
general theorems used without proof.
\end{abstract}
\tableofcontents
\newpage

\section{How to read these notes}
\label{sec:reading}

We work over an algebraically closed field $k$ of characteristic zero; taking
$k=\mathbb C$ throughout is entirely sufficient. Schemes and group actions in the
examples are of finite type. The group $\Gm$ is $\Spec k[t,t^{-1}]$, and
$\mu_n\subset\Gm$ is the group of $n$th roots of unity. Here $n\geq2$ unless
otherwise stated.

You should be comfortable with affine schemes, fibre products, sheaves, line
bundles, divisors, and $\Proj$. No knowledge of Hartshorne, Chapter III, is assumed.
Some ingredients that are not developed systematically in Chapters I--II---
\emph{étale morphisms}, \emph{torsors}, and \emph{algebraic spaces}---are introduced
when they are first needed. Their general theory is unnecessary for the worked
examples. Higher categories, derived geometry, intersection theory on stacks, and
the construction of moduli of curves are not prerequisites.

The main route is
\[
\begin{gathered}
\text{quotient rings}\longrightarrow\text{objects and their isomorphisms}
\longrightarrow\text{families and gluing}\\
\longrightarrow\text{quotient stacks}\longrightarrow\text{DM stacks}
\longrightarrow\text{line bundles and cuts}.
\end{gathered}
\]
Sections~\ref{sec:quotients}--\ref{sec:dm} give the foundations.
Sections~\ref{sec:coarse}--\ref{sec:roots} work out the geometry and sheaves.
Sections~\ref{sec:cutting}--\ref{sec:toycut} connect them to semiprojective
algebraic cutting. The exercises in Section~\ref{sec:exercises} have solutions.
An alphabetical glossary of terms, definitions, and notation is provided
in Appendix~\ref{sec:glossary}. Its companion Examples Glossary,
Appendix~\ref{sec:examples-glossary}, gives simple examples for every entry.

The notes use calligraphic letters for stacks and roman letters for their coarse
spaces. The notation $[X/G]$ always denotes a \emph{quotient stack}; an affine
invariant quotient is written $\Spec A^G$. These are different constructions.

\paragraph{A useful question to keep asking.}
When replacing an object by its isomorphism class, what have we forgotten about
its automorphisms, its families, or its line bundles? Each example below supplies
a concrete answer.

\section{What an invariant quotient forgets}
\label{sec:quotients}

\subsection{A quotient whose coarse space is smooth}
Let $\mu_n$ act on $\A^1_z$ by $z\mapsto\zeta z$. A polynomial is invariant
exactly when its monomials have degrees divisible by $n$, so
\[
 k[z]^{\mu_n}=k[z^n]=k[u],\qquad u=z^n.
\]
The usual affine quotient is another affine line,
\[
 q:\A^1_z\longrightarrow\A^1_u,\qquad z\longmapsto z^n.
\]
Over a nonzero geometric point $u$, the $n$ roots form one orbit and the action
on that orbit is free. At $z=0$, the orbit contains one point, but every
$\zeta\in\mu_n$ fixes it. Thus
\[
 \Stab(z)=\begin{cases}\{1\},&z\ne0,\\ \mu_n,&z=0.\end{cases}
\]
The coordinate ring of the quotient records invariant functions. It does not,
by itself, record this jump in stabilisers. In particular, smoothness of the
quotient $\A^1_u$ does not imply that the action was free.

One might try to decorate $u=0$ with the label $\mu_n$. That is a useful first
picture, but it is incomplete: we also need to say how the group acts on nearby
directions and on fibres of sheaves, and how such objects vary in families.

\subsection{A quotient whose coarse space is singular}
Let $\mu_2$ act on $\A^2_{x,y}$ by
\[
 (x,y)\longmapsto(-x,-y).
\]
An invariant monomial has even total degree. Hence
\[
 k[x,y]^{\mu_2}=k[a,b,c]/(ac-b^2),\qquad
 a=x^2,\quad b=xy,\quad c=y^2.
\]
For example, every even monomial is a product of $x^2,xy,y^2$, and the displayed
relation gives the familiar quadratic cone. The derivatives of $ac-b^2$ are
$c,-2b,a$, all zero at the vertex. The quotient is singular there, although
$\A^2$ is smooth.

We will obtain a smooth stack $[\A^2/\mu_2]$ with this cone as coarse space.
Its origin retains the sign action on the two-dimensional tangent space. The
previous example will give a smooth stack with a smooth coarse space. Taken
together, the examples show that the singularities of the coarse space do not
provide a complete test for the presence of stack structure.

\subsection{Retain arrows as well as orbits}
For a group $G$ acting on a set $S$, form a category with objects $x\in S$ and
arrows
\[
 \Hom(x,x')=\{g\in G:g\cdot x=x'\}.
\]
Composition is multiplication in $G$. Every arrow is invertible; a category with
this property is a \emph{groupoid}. The isomorphism classes are the orbits, while
$\Aut(x)=\Stab_G(x)$. Thus a groupoid keeps both the quotient points and their
symmetries.

There is an important scheme-level distinction. For the sign action on $\A^1$,
the action groupoid has arrow scheme $\mu_2\times\A^1$, which is two disjoint
affine lines. The relation obtained by comparing images in the coarse quotient
has scheme
\[
 \A^1\times_{\A^1_u}\A^1
 =\Spec k[z,z']/(z^2-z'^2).
\]
This is two lines meeting at the origin. The action groupoid has two distinct
arrows from $0$ to itself: the identity and the sign. The coarse relation has
lost their distinction. Replacing a groupoid by a relation is already a loss of
scheme-theoretic information.

\reading{Compute the stabilisers in both examples. Why could neither coarse
coordinate ring tell you which character occurs on an extra line over the origin?}

\section{From points to families: why a stack is groupoid-valued}
\label{sec:families}

\subsection{The functor-of-points viewpoint}
A scheme $X$ determines a contravariant functor
\[
 T\longmapsto X(T)=\Hom(T,X).
\]
A morphism $T'\to T$ pulls a $T$-point back to a $T'$-point. The Yoneda principle
says that the functor, with its behaviour on morphisms, determines $X$. One must
test on general schemes $T$, including nonreduced schemes, rather than just on
$\Spec k$.

For a moduli problem, a $T$-point is usually a \emph{family over $T$}. Two
families can be isomorphic in more than one way. The natural value of the moduli
problem is therefore a groupoid, not merely a set of isomorphism classes.

For example, let $\mathcal B(T)$ be the groupoid of line bundles on $T$, with
line-bundle isomorphisms as arrows. A map $f:T'\to T$ gives pullback
$L\mapsto f^*L$. Over $\Spec k$, all line bundles are isomorphic, but the
automorphism group of a one-dimensional vector space is $k^*$. Over $\PP^1$,
the objects include all $\OO_{\PP^1}(d)$. A single isomorphism class over a point
does not make this moduli problem the one-point scheme.

\subsection{Gluing must retain the chosen isomorphisms}
Recall how line bundles glue. For an open cover $\{T_i\}$ of $T$, choose line
bundles $L_i$ and isomorphisms
\[
 \alpha_{ij}:L_j|_{T_i\cap T_j}\xrightarrow{\sim}
                 L_i|_{T_i\cap T_j}.
\]
On triple overlaps require
\[
 \alpha_{ij}\alpha_{jk}=\alpha_{ik}.
\]
Then the $L_i$ glue to a line bundle on $T$, unique up to an isomorphism
compatible with the gluing data. Isomorphisms between two glued bundles also glue.
This is the model for the two stack axioms:
\begin{enumerate}
\item isomorphisms form a sheaf;
\item objects equipped with compatible gluing isomorphisms can be glued.
\end{enumerate}
The second assertion is called \emph{effective descent for objects}. Uniqueness
means uniqueness up to compatible isomorphism; it does not mean that all
automorphism groups vanish.

\begin{definition}[Stack, in working language]
A stack in groupoids assigns to every test scheme $T$ a groupoid $\cX(T)$,
with pullback functors for maps of test schemes, such that isomorphisms and
objects satisfy the two descent conditions above for the chosen topology.
Pullback along a composite agrees with successive pullback through specified
canonical isomorphisms satisfying a compatibility condition.
\end{definition}

The last sentence is the reason textbooks use \emph{categories fibred in
groupoids}. It formalises the fact that $(fg)^*L$ and $g^*f^*L$ are canonically
isomorphic, rather than insisting that they be literally the same set-theoretic
object. No further category theory will be needed for our calculations.

\subsection{Why étale covers enter}
An étale morphism is the algebraic analogue of a map that is locally an
isomorphism without introducing ramification. Formally, it is a flat,
unramified morphism locally of finite presentation. Open immersions are étale,
but there are many other examples.

A model calculation is
\[
 \Spec A[s]/(s^n-a)\longrightarrow\Spec A.
\]
If $a$ is a unit and $n$ is invertible in $A$, then $s$ is a unit upstairs and
the derivative $ns^{n-1}$ is invertible. This is a finite étale cover. It adjoins
an $n$th root of a unit, which generally cannot be chosen on any Zariski open
cover. An étale cover is a jointly surjective family of étale morphisms;
overlaps are fibre products $T_i\times_T T_j$.

In a full definition we can use the fppf topology: covers are jointly
surjective flat morphisms locally of finite presentation. It includes étale
covers. The stacks in these notes satisfy fppf descent. For our smooth groups,
their torsors can be trivialised étale locally, so thinking in terms of étale
covers suffices for the examples. The standard definitions of algebraic and DM
stacks used below are the fppf versions.

\subsection{The simplest twisted family}
Set $T=\Spec k[t,t^{-1}]$ and
\[
 P=\Spec k[s,s^{-1}],\qquad t=s^n.
\]
The group $\mu_n$ acts by $s\mapsto\zeta s$. Every geometric fibre is a free
transitive $\mu_n$-set. This is a \emph{$\mu_n$-torsor}. It is not the trivial
torsor: a section would give an $n$th root of $t$ in $k[t,t^{-1}]$. There is
not even such a root in $k(t)$, as the valuation at $t=0$ would give
$n\operatorname{ord}_0(s)=1$. The same argument excludes a section on any
nonempty Zariski open subset of $T$.

Nevertheless, after the étale base change $P\to T$, the torsor becomes trivial.
This is a concrete reason to enlarge the topology beyond Zariski covers.
It also demonstrates that algebraic closedness of $k$ does not trivialise
torsors over a parameter scheme.

Consider the presheaf sending $T$ to \emph{isomorphism classes} of
$\mu_n$-torsors. The torsor $P$ and the trivial torsor are distinct globally,
but become isomorphic on the cover $P\to T$. Hence that presheaf is not even a
sheaf: distinct sections can become equal on a cover. The groupoid remembers
the gluing isomorphisms that distinguish these families.

\source{For descent and the stack axioms, see Casalaina-Martin--Wise
\cite[\S\S2--3, especially Definition 3.31]{CMW} and Vistoli \cite[Chapters 3--4]{Vistoli}.}

\section{Quotient stacks and classifying stacks}
\label{sec:quotientstacks}

\subsection{A torsor is a family of copies of a group}
A left $G$-torsor $P\to T$ is a scheme with a left $G$-action which is locally
isomorphic, in the chosen topology, to $G\times T$ with translation on $G$.
Equivalently, it is locally nonempty and the map
\[
 G\times P\longrightarrow P\times_T P,\qquad(g,p)\longmapsto(gp,p)
\]
is an isomorphism. In the cases at hand the torsor map is faithfully flat and
locally of finite presentation. A torsor has no preferred identity section;
choosing a section trivialises it.

For $G=\Gm$, torsors are equivalent to line bundles: remove the zero section
from a line bundle, or form the associated line from a torsor. Thus torsors are
already familiar objects with the choice of a local frame left unspecified.

\subsection{Definition of the quotient stack}
\begin{definition}
Suppose $G$ acts on a scheme $X$ on the left. An object of $[X/G](T)$ is a
left $G$-torsor $P\to T$ together with a $G$-equivariant map
\[
 f:P\longrightarrow X,\qquad f(gp)=g f(p).
\]
An arrow is a $G$-equivariant torsor isomorphism commuting with the maps to $X$.
Pullback is obtained by base change.
\end{definition}

One can instead use right torsors and require $f(pg)=g^{-1}f(p)$. Both
conventions give the same stack. We use left torsors here so that all displayed
actions have the same direction.

Why include the torsor? After trivialising it on $T_i$, an object gives maps
$x_i:T_i\to X$. Different trivialisations are related by maps
$g_{ij}:T_i\times_T T_j\to G$, with
\[
 x_i=g_{ij}x_j,\qquad g_{ij}g_{jk}=g_{ik}.
\]
The torsor packages exactly these transition functions. Retaining only global
maps $T\to X$ modulo $G(T)$ would miss the objects with nontrivial transition
functions. The quotient stack is what results from allowing these local
presentations to glue.

Over an algebraically closed field, the torsors for the groups used here are
trivial, so the geometric-point groupoid really is the action groupoid:
\[
 \Aut_{[X/G]}(x)=\Stab_G(x).
\]
This formula is an excellent first calculation, but is not a substitute for
the definition over general $T$.

\begin{example}[The classifying stack]
Take $X=\Spec k$ with trivial $G$-action. Then
\[
 BG=[\Spec k/G],\qquad BG(T)=\{G\text{-torsors over }T\}.
\]
For $G=\mu_n$, there is one geometric isomorphism class, with automorphism
group $\mu_n$. The coarse space will be a point. A family over $T$ still
remembers its torsor, including the nontrivial example of the preceding section.

Equivalently, $B\mu_n(T)$ is the groupoid of pairs
$(N,\phi:N^{\otimes n}\xrightarrow{\sim}\OO_T)$.
The isomorphism $\phi$ is part of the data. Keeping just a line bundle whose
$n$th power happens to be trivial would lose information.
\end{example}

\begin{example}[A free action]
If $X\to Y$ is a $G$-torsor, then $[X/G]\simeq Y$. Indeed, a map $T\to Y$
pulls $X$ back to a torsor over $T$ with an equivariant map to $X$. Conversely,
such a pair descends to a map to $Y$. Thus stacks include ordinary free
quotients. For example,
\[
 [\Gm/\mu_n]\simeq\Gm,\qquad z\longmapsto z^n.
\]
\end{example}

\subsection{Atlases and the first stack fibre product}
There is a natural map $p:X\to[X/G]$, given by the trivial torsor. If
$T\to[X/G]$ corresponds to $(P,f)$, then
\begin{equation}\label{eq:atlasbasechange}
 T\times_{[X/G]}X\simeq P.
\end{equation}
The extra data in this fibre product is a trivialisation of the pulled-back
object, equivalently a point of its torsor. This explains why $p$ is
surjective locally even when an object over $T$ has no global lift to $X$.

If $G$ is smooth, $p$ is a smooth cover. If $G$ is finite étale, $p$ is an
étale cover. These assertions follow directly from
\eqref{eq:atlasbasechange}, since they hold for $P\to T$.

In particular, $\A^1\to[\A^1/\mu_n]$ is étale \emph{at the origin as well}.
Its composite with the coarse map is $z\mapsto z^n$, which is ramified there.
The atlas map and the coarse map have different targets, so there is no
contradiction.

A fibre product of stacks is a \emph{2-fibre product}: an object of
$\cX\times_{\cY}\cZ$ consists of $x$, $z$, and a specified isomorphism between
their images in $\cY$. This is the only extra ingredient beyond the fibre
products of schemes. For the atlas above it gives
\[
 X\times_{[X/G]}X\simeq G\times X.
\]
The two projections are source and target, $x$ and $gx$. The arrows, their
composition, and the object scheme form a \emph{groupoid presentation} of
the stack.
The atlas calculation is a 2-Cartesian square:
\[
\begin{tikzcd}
 P \arrow[r,"f"] \arrow[d] & X \arrow[d,"p"]\\
 T \arrow[r,"{(P,f)}"'] & {[X/G]}.
\end{tikzcd}
\]
The square includes the isomorphism identifying the two induced objects
on $P$. For $T=X$, the square recovers the arrow space $G\times X$.

\source{The torsor construction and its equivalence with the quotient-groupoid
construction are given in the Stacks Project, \href{https://stacks.math.columbia.edu/tag/04UV}{Tag 04UV}.
The algebraicity of smooth groupoid quotients is \href{https://stacks.math.columbia.edu/tag/04TK}{Tag 04TK}.}

\section{The definition of a Deligne--Mumford stack}
\label{sec:dm}

\subsection{A short bridge through algebraic spaces}
An \emph{algebraic space} is a sheaf of sets on the étale site which has a
representable diagonal and an étale surjective cover by a scheme. Here
representability of the diagonal is by schemes. An algebraic space is thus
étale locally a scheme, with ordinary set-valued points rather than groupoid-valued
points. Every scheme gives an algebraic space. The converse is false in general.

This enlargement lets quotients exist in situations where Zariski gluing of
schemes is too restrictive. We do not need an explicit nonschematic algebraic
space in these notes: every displayed coarse space is a scheme. The phrase
``algebraic space'' in the formal definitions prevents an unnecessary
restriction on the general theory.

\subsection{Representability and the diagonal}
A morphism $\cX\to\cY$ is \emph{representable by algebraic spaces} if, after
every map $T\to\cY$ from a scheme, the fibre product $T\times_{\cY}\cX$ is
an algebraic space. A representable morphism is called étale, smooth, finite,
and so on when its base changes to schemes have that property.

The diagonal
\[
 \Delta_{\cX}:\cX\longrightarrow\cX\times\cX
\]
organises isomorphisms. Given two objects $\xi,\eta\in\cX(T)$, its pullback
along $(\xi,\eta):T\to\cX\times\cX$ is
\[
 \underline{\Isom}_T(\xi,\eta).
\]
Over a map $T'\to T$, this sheaf gives isomorphisms between the pulled-back
families. When $\xi=\eta$, it is the automorphism group space of the family.
Thus representability of the diagonal means that isomorphisms between
families themselves form algebraic spaces.

For quotient objects represented by maps $x,x':T\to X$, this is the
transporter
\[
 \{g\in G_T:g x=x'\}.
\]
For the sign action at $x=x'=0$, it is $\mu_2\times T$. For a scheme
considered as a stack, there are only identity arrows: the diagonal expresses
equality, rather than a nontrivial space of isomorphisms.

\begin{definition}[Algebraic stack and DM stack]
An \emph{algebraic stack} is an fppf stack in groupoids whose diagonal is
representable by algebraic spaces and which admits a representable smooth
surjective morphism $U\to\cX$ from a scheme. Such a map is called a smooth atlas.

An algebraic stack is \emph{Deligne--Mumford}, or \emph{DM}, if it admits a
representable étale surjective morphism $U\to\cX$ from a scheme. Such a map
is an étale atlas.
\end{definition}

An étale atlas provides scheme charts with no continuous choices remaining
in each fibre. The stack can still have finite automorphism groups because
different chart identifications are retained as arrows.

We follow the Stacks Project definition. Some books build quasi-compactness
or separatedness conditions on the diagonal into the terminology. Here those
conditions are stated separately. In applications to coarse spaces and descent
below, we impose finite type and separatedness.

\begin{theorem}[Standard criterion; used without proof]\label{thm:dmcriterion}
An algebraic stack is DM if and only if its diagonal is unramified.
In the finite-type quotient situations considered here, over characteristic
zero, $[X/G]$ for a smooth algebraic group $G$ is DM if and only if every
geometric stabiliser is finite.
\end{theorem}

The diagonal formulation says that automorphisms have no infinitesimal motion.
In our finite-type quotient examples, or for the separated finite-type
stacks used below, the stabilisers are finite and reduced. With the broad
Stacks Project convention, DM alone does not impose quasi-compactness
of the diagonal; finiteness requires the additional hypotheses just stated. Over
characteristic zero, finite subgroup schemes of our algebraic groups are
reduced. Outside characteristic zero one must retain the reduced/unramified
condition: a finite group scheme such as $\mu_p$ in characteristic $p$ is
not enough.

\subsection{Examples, nonexamples, and what the definition does not assert}
\begin{enumerate}
\item Every scheme and every algebraic space is a DM stack, with discrete
groupoids of points. An algebraic stack with trivial geometric stabilisers is
an algebraic space; it need not be a scheme.
\item For a finite group $H$ over $k$, the map $X\to[X/H]$ is an étale atlas.
Thus $B\mu_n$, $[\A^1/\mu_n]$, and $[\A^2/\mu_2]$ are DM. The last two
are smooth because they have smooth étale atlases.
\item $B\Gm$ is algebraic, but not DM. Its geometric object has automorphism
group $\Gm$. It is the stack of line bundles introduced earlier.
\item $[\A^1/\Gm]$ for the scaling action is not DM: the origin has stabiliser
$\Gm$, although the nonzero locus has trivial stabiliser.
\item A quotient by a positive-dimensional group can be DM. For example,
$[(\A^2\setminus0)/\Gm]$ with positive weights $(n,1)$ has only finite
stabilisers. Its obvious atlas is smooth of relative dimension one, rather
than étale. Other étale charts exist; we construct them in
Section~\ref{sec:weighted}.
\end{enumerate}

The DM condition does not imply smoothness: take any singular scheme. It
does not imply separatedness: take any nonseparated scheme. Nor does it assert
that a given quotient presentation is by a finite group. It is a local
condition on a stack with an algebraic presentation.

\paragraph{Separatedness.}
A stack is separated over $k$ if its diagonal is proper. For a finite-type
DM stack the diagonal is already unramified and locally quasi-finite; if it is proper
it is finite. This explains why finite stabilisers and well-behaved finite
quotient charts occur together in separated DM geometry. Closed orbits and
separation questions must still be checked in a GIT application; finite
stabilisers alone are not a general substitute for those checks.

\paragraph{Dimension and smoothness.}
For a smooth group $G$ acting on $X$, the smooth atlas has relative dimension
$\dim G$, so $\dim[X/G]=\dim X-\dim G$ in the equidimensional examples.
The quotient stack is smooth over $k$ when $X$ is smooth: smoothness can be
checked on a smooth atlas. Thus $\PP(n,1)$ is a smooth curve, while
$[\A^2/\mu_2]$ is a smooth surface.

\paragraph{Relation to orbifolds.}
Over $\mathbb C$, a smooth separated DM stack behaves locally like a complex
manifold modulo a finite group. This is the useful orbifold picture. Some
definitions of orbifold require every local group to act faithfully. DM
stacks allow ineffective actions: $B\mu_n$ is a simple example, since its
tangent space is zero. A singular DM stack is also outside the smooth
orbifold picture. We will therefore use the algebraic definition rather than
identify the two terms without qualifications.

\source{Definitions: Stacks Project, \href{https://stacks.math.columbia.edu/tag/026N}{Tag 026N}.
The diagonal criterion is \href{https://stacks.math.columbia.edu/tag/06N3}{Tag 06N3}.
For the diagonal as an isomorphism sheaf, see \cite[Example 4.6]{CMW}.}

\section{Coarse spaces and local finite quotients}
\label{sec:coarse}

\begin{definition}[Coarse moduli space]
A morphism $\pi:\cX\to X$ to an algebraic space is a coarse moduli space if:
\begin{enumerate}
\item for every algebraically closed extension $\Omega/k$, it induces a
bijection between isomorphism classes in $\cX(\Omega)$ and $X(\Omega)$;
\item it is universal for morphisms from $\cX$ to algebraic spaces: every
such morphism factors uniquely through $X$.
\end{enumerate}
\end{definition}

Both clauses matter. The first says that geometric isomorphism classes are
preserved as points. The second controls morphisms and functions. Neither
says that families over general schemes are preserved.

For a finite group $H$ acting on an affine scheme $\Spec A$, the coarse space
of $[\Spec A/H]$ is $\Spec A^H$. A morphism from the stack to a scheme is
the same as an $H$-invariant morphism from the atlas, so the invariant-ring
calculation is consistent with the universal property. The universality for
general algebraic-space targets is part of the quotient theorem.

\begin{theorem}[Coarse spaces and local structure; used without proof]
A separated finite-type DM stack over $k$ admits a coarse moduli space,
which is an algebraic space. Étale locally on this coarse space it has
finite-group quotient presentations. At a chosen geometric point, a chart
can be chosen with the stabiliser of that point as the local finite group.
\end{theorem}

This gives a way to learn the general theory through finite quotient
calculations. It is not a statement that every stack has one global
presentation by a finite group, nor that its coarse space must be a scheme.

\subsection{A coarse map preserves points but changes families}
The coarse space of $B\mu_n$ is $\Spec k$. There is one geometric
isomorphism class on each side. But a map $T\to B\mu_n$ carries a torsor,
whereas a map $T\to\Spec k$ carries no such data. The nontrivial torsor
$s^n=t$ and the trivial torsor are sent to the same map to the coarse point.

The coarse space of $[\A^1/\mu_n]$ is $\A^1_u$. Its nontrivial stabiliser
is invisible in that smooth line. In the quadratic-cone example, the coarse
space instead has a singularity. Thus the implication
``smooth stack $\Rightarrow$ smooth coarse space'' is false, and
``smooth coarse space $\Rightarrow$ no stack structure'' is also false.

\subsection{Good moduli spaces are a different notion}
For a reductive GIT quotient with strictly semistable points, different
orbits can have closures meeting the same closed orbit. The GIT quotient
identifies them. Such a quotient need not be a coarse moduli space of the
whole semistable quotient stack.

The smallest example is $[\A^1/\Gm]$ for scaling. It has two geometric
isomorphism classes: the zero orbit and the nonzero orbit. Its affine
invariant quotient is $\Spec k$, since $k[z]^{\Gm}=k$. That point is
\emph{not} its coarse moduli space: it fails the bijection in the definition.

This quotient is a \emph{good moduli space}. For our purposes its defining
properties are that the pushforward on quasi-coherent sheaves is exact and
$\OO_X\simeq\pi_*\OO_{\cX}$; the morphism is also required to be
quasi-compact. These conditions extend the behaviour of invariant rings and
invariant modules to stacks. We will use the term only to distinguish the
quotients at critical cutting levels from the coarse quotients at regular ones.

\source{Coarse-space existence is discussed in the Stacks Project,
\href{https://stacks.math.columbia.edu/tag/04UX}{Tag 04UX}; local structure is
discussed in \cite[\S3.4]{Behrend}. For good moduli spaces see
\cite[Definition 4.1 and \S13]{Alper}.}

\section{Sheaves and line bundles retain stabiliser representations}
\label{sec:sheaves}

\subsection{Equivariant sheaves are sheaves on the quotient stack}
A vector bundle on $[X/G]$ is equivalent to a $G$-equivariant vector bundle
on $X$. An equivariant structure says how $g$ identifies the fibre at $x$
with the fibre at $gx$, compatibly with multiplication in $G$. This is the
descent datum on the atlas and its arrow space $G\times X$.

At a point fixed by $g$, the equivariant structure is an action on the same
fibre. Hence a vector bundle on a stack gives a representation of every
stabiliser. In rank one this representation is a \emph{character}.
These characters are additional geometric data, even if the underlying
bundle on the atlas is trivial.

\begin{example}[All line bundles on $B\mu_n$]
A vector bundle on $B\mu_n$ is a representation of $\mu_n$. Its line bundles
are the one-dimensional representations
\[
 \chi^c:\mu_n\to\Gm,\qquad\zeta\longmapsto\zeta^c,
 \qquad c\in\mathbb Z/n\mathbb Z.
\]
Thus $\Pic(B\mu_n)\simeq\mathbb Z/n\mathbb Z$, while the coarse point has
trivial Picard group. The global sections of the character line are its
invariant vectors: $k$ for $c=0$, and zero for $c\ne0$.
\end{example}

\subsection{Character conventions and invariant sections}
On $[\A^1_z/\mu_n]$, let $\cL_c$ be the trivial line upstairs with the
\emph{geometric fibre action}
\[
 (z,v)\longmapsto(\zeta z,\zeta^c v).
\]
A section descends to the stack if it is equivariant. Writing it as the
function $f(z)$ in this trivialisation gives
\begin{equation}\label{eq:equivsection}
 f(\zeta z)=\zeta^c f(z).
\end{equation}
Thus the allowable monomials are $z^j$ with $j\equiv c\pmod n$.

This equation avoids an ambiguity about representations on functions. With
the inverse-pullback convention $g\cdot f=f\circ g^{-1}$, the coordinate $z$
has weight $-1$. With the direct-pullback convention $f\mapsto f\circ g$ for
our abelian groups, it has weight $+1$. The \emph{geometric tangent action}
$z\mapsto\zeta z$ is fixed throughout. Always state which action defines a
weight before comparing formulas.

\subsection{Pushforward is not descent}
Let $\pi:[\A^1/\mu_n]\to\A^1_u$, $u=z^n$. Choose $r$ with
$0\leq r<n$ and $r\equiv c\pmod n$. Equation~\eqref{eq:equivsection} gives
\begin{equation}\label{eq:rankonepush}
 \Gamma(\A^1_u,\pi_*\cL_c)=z^r k[u].
\end{equation}
As a $k[u]$-module this is free of rank one. Therefore the pushforward is an
ordinary line bundle, even when $c\not\equiv0\pmod n$.

Nevertheless $\cL_c$ does not descend from the coarse space when
$c\not\equiv0\pmod n$.
If $\cL_c\simeq\pi^*L$, the stabiliser at $z=0$ would act trivially on its
fibre: the coarse point has no nontrivial automorphisms that could act on $L$.
But that fibre has character $c$. The natural comparison
\[
 \pi^*\pi_*\cL_c\longrightarrow\cL_c
\]
is locally given by multiplication by $z^r$ in the displayed trivialisation.
For $r>0$ it vanishes at $0$, so it is not an isomorphism.

There are therefore three different operations or questions:
\begin{enumerate}
\item the original line bundle $\cL$ on the stack;
\item its sheaf $\pi_*\cL$ on the coarse space;
\item whether $\cL$ is isomorphic to the pullback of a line bundle there.
\end{enumerate}
Even an invertible pushforward does not answer the third question.

\begin{theorem}[Descent criterion in our setting; used without proof]
Let $\pi:\cX\to X$ be the coarse map of a separated finite-type DM stack
over $k$. A vector bundle $E$ on $\cX$ is pulled back from a vector bundle
on $X$ if and only if every geometric stabiliser acts trivially on its fibre.
In particular this applies to line bundles.
\end{theorem}

Characteristic zero makes the finite stabilisers linearly reductive, so these
coarse maps are tame good moduli spaces. The necessity of the criterion was
proved above. The sufficiency is a general descent theorem, rather than a
consequence of counting sections. In a finite quotient presentation it says
that equivariant gluing with trivial fibre actions is ordinary quotient
gluing. See Alper's vector-bundle descent theorem \cite[Theorem 10.3]{Alper}.

\subsection{An invariant module which is not a line bundle}
Return to $\cX=[\A^2_{x,y}/\mu_2]$ with simultaneous sign action, and give a
trivial line upstairs the sign fibre character. Set
\[
 A=k[x^2,xy,y^2],\qquad
 M=\{f\in k[x,y]:f(-x,-y)=-f(x,y)\}.
\]
The pushforward is the coherent sheaf associated to the $A$-module $M$.
It is generated by $x,y$. At the cone vertex $\mathfrak m=(x^2,xy,y^2)$,
the vector space $M/\mathfrak m M$ has dimension two: the classes of $x,y$
remain independent. On the other hand $M$ has rank one over the fraction
field of $A$, because $y=(xy/x^2)x$ generically. A free rank-one module at
the vertex would have a one-dimensional quotient modulo $\mathfrak m$.
Thus $M$ is not locally free there.

This is a concrete model for why the pushforward of an honest stack line
bundle can fail to be invertible on the coarse space. It complements
\eqref{eq:rankonepush}, where the pushforward is invertible but descent
still fails.

\paragraph{Why invariants behave well here.}
For a finite group $H$ in characteristic zero, averaging
$v\mapsto |H|^{-1}\sum_{h\in H}h v$ projects onto invariant vectors.
Consequently taking invariants is exact. On affine finite-quotient charts,
pushforward is the corresponding invariant-module operation. Exactness is
very useful, but it does not make pushforward preserve tensor products, as
the next example will show.

\source{Equivariant sheaves and invariant sections are explained in
\cite[\S3.5, especially Exercise 3.56]{Behrend}.}

\section{A weighted projective line, worked out in charts}
\label{sec:weighted}

Let $\Gm$ act on $U=\A^2_{x,y}\setminus\{0\}$ by
\begin{equation}\label{eq:weightedaction}
 t\cdot(x,y)=(t^n x,t y).
\end{equation}
Define
\[
 \cP=\PP(n,1):=[U/\Gm].
\]
This is the same weighted projective line often written $\PP(1,n)$ after
interchanging the coordinates. We retain the order $(x,y)$ to match the cut
in Section~\ref{sec:toycut}.

If $y\ne0$, a stabilising $t$ satisfies $t=1$. At $(x,0)$ with $x\ne0$,
it satisfies $t^n=1$. All stabilisers are finite, so $\cP$ is DM. It is
smooth because $U$ is smooth.

\subsection{The two charts}
On $y\ne0$, there is a unique way to set $y=1$: choose $t=y^{-1}$. The
remaining coordinate is
\[
 z=\frac{x}{y^n}.
\]
Thus this open substack is an ordinary affine line $\A^1_z$.

On $x\ne0$, setting $x=1$ requires a choice $t^n=x^{-1}$. Such a choice
exists étale locally, but has a $\mu_n$ ambiguity. The remaining coordinate
$s$ transforms by $s\mapsto\zeta s$. Consequently this chart is
\[
 [\A^1_s/\mu_n].
\]
Its coarse coordinate is
\[
 u=s^n=\frac{y^n}{x}.
\]
On the overlap $zu=1$. The two coarse affine lines therefore glue to
$\PP^1$. The coarse map is
\begin{equation}\label{eq:coarseweighted}
 \pi:\cP\longrightarrow\PP^1,\qquad [x:y]\longmapsto[x:y^n].
\end{equation}
The expression on the right is unchanged as a projective point when
$(x,y)$ is replaced by $(t^n x,ty)$.

Taking the ordinary affine line atlas of the second quotient chart and the
first affine chart gives an étale atlas of $\cP$. This constructs the atlas
whose existence was asserted by the DM criterion. It is different from
the smooth $\Gm$-torsor atlas $U\to\cP$.

\subsection{A description in families}
An object of $\cP(T)$ can be described as
\[
 (N,\,x\in\Gamma(T,N^{\otimes n}),\,
        y\in\Gamma(T,N)),
\]
where the two sections do not vanish simultaneously at any geometric point.
An isomorphism of line bundles must carry the displayed sections to the
corresponding sections. This is the family version of weighted coordinates.
In the left-torsor convention above, the corresponding torsor is
$P=(N^\vee)^\times$: a nonzero covector evaluates the sections to give
the coordinates. Scaling that covector by $t$ multiplies those
coordinates by $t^n$ and $t$, as required.

For example, where $y$ is nonzero it trivialises $N$, leaving the ordinary
function $x/y^n$. Where $x$ is nonzero it trivialises $N^{\otimes n}$,
leaving an $n$th-root choice for $N$. Thus the chart calculation reflects
actual families, rather than merely orbit notation.

\subsection{The boundary and the line bundles}
Let $\cD=\{y=0\}\subset\cP$. It is
\[
 \cD=[\Gm/\Gm]\simeq B\mu_n,
\]
where the action on the first $\Gm$ is $t\cdot x=t^n x$. Its coarse space
is the point $D=[1:0]\in\PP^1$.

Define $\OO_{\cP}(c)$ using the geometric fibre action
\[
 (x,y,v)\longmapsto(t^n x,ty,t^c v).
\]
Its global sections are weighted homogeneous polynomials of degree $c$.
The coordinate $y$ is an equivariant section of $\OO_{\cP}(1)$ and vanishes
simply along $\cD$. Hence
\begin{equation}\label{eq:weightlines}
 \OO_{\cP}(1)\simeq\OO_{\cP}(\cD),\qquad
 \pi^*\OO_{\PP^1}(1)\simeq\OO_{\cP}(n).
\end{equation}
The second identity can also be read directly from the two degree-$n$
sections $x,y^n$ defining the coarse map.

At $\cD$, the stabiliser acts on $\OO_{\cP}(c)$ by $\zeta^c$.
It follows that
\begin{equation}\label{eq:weighteddescent}
 \OO_{\cP}(c)\text{ descends to }\PP^1
 \quad\Longleftrightarrow\quad n\mid c.
\end{equation}
Sufficiency follows already from \eqref{eq:weightlines}; necessity follows
from the fibre character. We have therefore verified this case of the
descent theorem without using its general proof.

\subsection{Counting sections and computing pushforward}
For $m\ge0$, the monomials of degree $m$ are
\[
 x^a y^{m-na},\qquad 0\leq a\leq\lfloor m/n\rfloor.
\]
Consequently
\begin{equation}\label{eq:weightedh0}
 \dim H^0(\cP,\OO_{\cP}(m))=\lfloor m/n\rfloor+1.
\end{equation}
No higher-cohomology machinery is needed: this just counts equivariant
polynomials on $U$. A regular function on $\A^2\setminus0$ extends to
$\A^2$; explicitly, $k[x,y,y^{-1}]\cap k[x,y,x^{-1}]=k[x,y]$ in $k(x,y)$.

More precisely, for every $m\in\mathbb Z$,
\begin{equation}\label{eq:weightedpush}
 \pi_*\OO_{\cP}(m\cD)
 \simeq\OO_{\PP^1}\bigl(\lfloor m/n\rfloor D\bigr).
\end{equation}
To see this locally near $D$, use a root coordinate $s$ with $u=s^n$.
The divisor sheaf upstairs is $s^{-m}k[s]$, and its invariant rational
monomials are $s^{nj}=u^j$ with $nj\ge-m$. Thus
\[
 (s^{-m}k[s])^{\mu_n}=u^{-\lfloor m/n\rfloor}k[u].
\]
Off $D$ both divisor sheaves are canonically trivial. This proves the formula
by gluing and includes negative $m$.

When $0<m<n$, the pushforward in \eqref{eq:weightedpush} is $\OO_{\PP^1}$,
but the stack line does not descend. In particular,
\[
 (\pi_*\OO_{\cP}(1))^{\otimes n}\simeq\OO_{\PP^1},\qquad
 \pi_*(\OO_{\cP}(1)^{\otimes n})\simeq\OO_{\PP^1}(1).
\]
Pushforward has forgotten stabiliser characters, and it does not commute
with this tensor product.

\reading{For $n=3$, write down the sections of $\OO(m)$ for $0\le m\le6$.
Which powers descend? Can the dimensions of sections alone recover the
character at the boundary?}

\section{Root stacks, divisors, and normal lines}
\label{sec:roots}

\subsection{Making the root of a divisor into geometric data}
Let $D\subset X$ be an effective Cartier divisor. Its line bundle $\OO_X(D)$
has a canonical section $s_D$ with zero divisor $D$. An $n$th root of this
pair should be a line bundle $N$ with a section $s$ and an isomorphism
$N^{\otimes n}\simeq\OO_X(D)$ carrying $s^n$ to $s_D$. Such a pair
need not exist on $X$. Rather than choose one, we parameterise all choices.

\begin{definition}[Divisor root stack]
An object of $\sqrt[n]{(X,D)}$ over a scheme $T$ is a map $f:T\to X$ and
data
\[
 (N,s,\phi),\qquad
 \phi:N^{\otimes n}\xrightarrow{\sim}f^*\OO_X(D),\qquad
 \phi(s^{\otimes n})=f^*s_D.
\]
Arrows are line-bundle isomorphisms compatible with $s$ and $\phi$.
\end{definition}

If $D$ has equation $a\in A$ on $\Spec A$, its local presentation is
\begin{equation}\label{eq:rootchart}
 \sqrt[n]{(\Spec A,D)}
 \simeq\left[\Spec A[z]/(z^n-a)\,/\mu_n\right],\qquad
 \zeta\cdot z=\zeta z.
\end{equation}
Changing a trivialisation changes the root coordinate and glues these
quotient charts. Since $n$ is invertible, this is a DM stack. If $X$ and
$D$ are smooth, it is smooth: locally $a$ is a regular coordinate and the
equation $a=z^n$ replaces it by the smooth coordinate $z$. For singular
$D$, smoothness is not automatic.

Away from $D$, the section $s$ is nonzero and trivialises $N$, so the root
stack is just $X\setminus D$. Over $D$, scaling the root by $\zeta$ preserves
its $n$th power. The stabiliser is $\mu_n$. Locally the invariant ring in
\eqref{eq:rootchart} is $A$, so the coarse space is $X$.

Our weighted line is the example
\[
 \PP(n,1)\simeq\sqrt[n]{(\PP^1,D)},\qquad D=[1:0].
\]
Its two charts are exactly the ordinary chart away from $D$ and
$[\A^1_s/\mu_n]$ near $D$.

\subsection{The reduced root divisor and the scheme-theoretic fibre}
The universal root section vanishes along a Cartier substack $\cD$ with
\[
 \pi^*D=n\cD,\qquad
 \OO(n\cD)\simeq\pi^*\OO_X(D).
\]
If $D$ is smooth, this $\cD$ is the reduced substack over $D$. Its local
description is $[D/\mu_n]$ with trivial action on $D$. Globally it can be
a nontrivial $\mu_n$-gerbe: an object over $T\to D$ is an $n$th root of
$\OO_X(D)|_T$. This is locally $D\times B\mu_n$, but a global product
would require a compatible global root. When $D$ is the single $k$-point
in our weighted-line example, $\cD\simeq B\mu_n$.

There is a useful scheme-theoretic subtlety. For $[\A^1_z/\mu_n]\to\A^1_u$
with $u=z^n$, the actual fibre over $u=0$ is
\[
 \left[\Spec k[z]/(z^n)\,/\mu_n\right].
\]
It is a nonreduced thickening of $B\mu_n$, rather than just $B\mu_n$.
The reduced residual gerbe, cut out by $z=0$, is $B\mu_n$.
The distinction is invisible on algebraically closed field-valued points
but visible on nonreduced test schemes. It is another reason to retain
the full functor of points.

\subsection{Normal and conormal characters}
At the origin of the chart $\A^1_z$, a normal vector transforms as
$v\mapsto\zeta v$. The normal line of $B\mu_n$ in $[\A^1/\mu_n]$
therefore has character $+1$. The conormal line is its dual and has
character $-1$. This statement concerns the geometric action on vectors
and covectors, independently of which convention is used to write an
action on functions.

If $i:\cD\hookrightarrow\cX$ is a Cartier divisor, its normal line is
\[
 \normal_{\cD/\cX}\simeq\OO_{\cX}(\cD)|_{\cD}.
\]
For the weighted line this is the character-$1$ line on $B\mu_n$.
The coarse divisor $D\subset\PP^1$ also has an ordinary normal line.
Its pullback to $B\mu_n$ has trivial stabiliser character, so it is not
the stack normal line. Instead,
\begin{equation}\label{eq:normalroot}
 \normal_{\cD/\cX}^{\otimes n}
 \simeq (\pi|_{\cD})^*\normal_{D/X}.
\end{equation}
This follows by restricting $\OO(n\cD)=\pi^*\OO(D)$.
Thus even when both spaces and their coarse divisors are smooth, the
normal line on the stack can carry information lost by the coarse pair.

The line $\OO(\cD)$ is sometimes thought of as the line of the
``fractional divisor'' $D/n$. It is an honest line bundle on the stack.
The fractional notation on the coarse space is a way of remembering the
root, not an assertion that $D/n$ is an integral Cartier divisor there.

\paragraph{Root stacks versus root gerbes.}
If one asks only for $N^{\otimes n}\simeq L$ and does not include a
section, the resulting root \emph{gerbe} has $\mu_n$ stabiliser everywhere.
The divisor root stack includes a section and has trivial stabiliser off
its zero locus. For instance, the root gerbe of the trivial line on $X$
is $X\times B\mu_n$, whereas $\sqrt[n]{(\A^1,0)}$ has stabiliser only
at the origin. Here $0$ in the latter notation means the origin divisor.

\source{The divisor root construction is due to Cadman \cite[\S2]{Cadman}.
For the concrete curve description and divisor pushforward formula, see
\cite[\S3.6, Definition 3.62 and Theorem 3.64]{Behrend}.}

\section{Where these ideas enter algebraic cutting}
\label{sec:cutting}

\subsection{The setting and the chosen sign convention}
Suppose $M$ is a smooth variety with a $\Gm$-action and an ample
$\Gm$-linearised line bundle $L$. A linearisation specifies maps on fibres
$t_L:L_p\to L_{tp}$ compatible with the group action. The section ring is
\[
 R=\bigoplus_{m\ge0}H^0(M,L^{\otimes m}),\qquad
 R_{m,w}=H^0(M,L^{\otimes m})_w.
\]
The index $m$ is a tensor or projective degree, while $w$ is a group weight.
They are independent. For semiprojective examples the degree-zero ring
$R_0=H^0(M,\OO_M)$ can be nonconstant.

For this discussion assume $R$ is finitely generated, $M\simeq\Proj R$,
and the affinisation $M\to\Spec R_0$ is projective. A contracting action
is an additional condition, not a consequence of being projective over an
affine base. Our toy model $M=\A^1$ has both properties.

We use the convention in the repository's
\path{claude-notes/conventions.tex}, dated 12 November 2026. For sections,
\[
 (t\cdot s)(p)=t_L^{-1}s(tp).
\]
For functions this is direct pullback, $f(p)\mapsto f(tp)$. Because our
group is abelian, these operators form a representation. A fixed component
$F$ has section weight $\Phi(F)$, so the \emph{geometric fibre weight}
of $L|_F$ is $-\Phi(F)$.

Introduce an auxiliary coordinate $y$ with point action
$y\mapsto t^\epsilon y$, where $\epsilon\in\{+1,-1\}$. Twist by a trivial
line with geometric fibre weight $d\in\mathbb Z$, whose constant section
$\one_d$ has section weight $-d$. On
$X_\epsilon=M\times\A^1_y$, the section
\[
 r\,y^\ell\one_d^m,\qquad r\in R_{m,w},
\]
is invariant exactly when
\begin{equation}\label{eq:cutinvariance}
 w+\epsilon\ell-md=0,\qquad \ell\ge0.
\end{equation}
It follows that the invariant section algebra is
\begin{align}
 B^+_m&=\bigoplus_{w\le md}R_{m,w}y^{md-w}\one_d^m,\label{eq:cutplus}\\
 B^-_m&=\bigoplus_{w\ge md}R_{m,w}y^{w-md}\one_d^m.\label{eq:cutminus}
\end{align}
These equations, rather than an unstated moment-map sign, determine the
two directions in this note.

\subsection{Stack cut, coarse cut, and boundary}
Let $X^{\sst}_\epsilon(d)$ be the locus on which some invariant section
of positive tensor degree is nonzero. The \emph{stack cut} and its GIT
quotient are
\[
 \cM^{\leq d}=[X^{\sst}_+(d)/\Gm],\quad
 M^{\leq d}=\Proj B^+,
\]
and similarly $\cM^{\geq d}$ and $M^{\geq d}$ using $\epsilon=-1$.
At regular levels in the setting below, the GIT quotients are coarse
spaces. At a general level they are good moduli spaces.

Putting $y=0$ in \eqref{eq:cutinvariance} leaves $w=md$. The common
boundary is therefore described by
\[
 \cM_d=[M^{\sst}(d)/\Gm],\qquad
 M_d=\Proj\bigoplus_mR_{m,md}.
\]
The stack cut carries the line bundle $\cL$ defined by the product
linearisation. The corresponding graded sheaf $\widetilde{B^\epsilon(m)}$
on the quotient is its pushforward in this setting. It need not be
invertible, and invertibility alone would still not prove that the stack
line descends. Section~\ref{sec:sheaves} gives both kinds of example.

\subsection{Why the stabilisers are concentrated on the boundary}
A point $(p,y)$ is fixed by $t$ only if
\[
 tp=p,\qquad t^\epsilon y=y.
\]
If $y\ne0$ and $\epsilon=\pm1$, the second equation forces $t=1$.
If $y=0$, the stabiliser is $\Stab_{\Gm}(p)$. Thus
\begin{equation}\label{eq:cutstabs}
 \Stab(p,y)=
 \begin{cases}
 \{1\},&y\ne0,\\
 \Stab_{\Gm}(p),&y=0.
 \end{cases}
\end{equation}
Closed algebraic subgroups of $\Gm$ over $k$ are $\Gm$ itself or finite
cyclic groups $\mu_r$ (including the trivial group).

In the current repository setting, a level is called regular when
$d\ne\Phi(F)$ for all fixed components. Then the semistable locus has
finite stabilisers and closed orbits within that locus. Its quotient stack
is smooth DM and has a geometric coarse quotient. Regularity permits
finite cyclic stabilisers on the boundary; it does not mean that the action
is free. At a critical level, fixed points can remain semistable and yield
$\Gm$ stabilisers. The resulting quotient stack is then algebraic but
not DM along those points.

One can already see the critical equality from section weights. At a fixed
point $p\in F$, a nonvanishing section of $L^{\otimes m}$ has weight
$m\Phi(F)$. At $y=0$, invariance demands $m\Phi(F)=md$. Thus fixed
points can survive on the boundary exactly at the critical weight.

\subsection{Normal character versus polarisation character}
On the regular boundary the auxiliary normal direction transforms by
$t^\epsilon$. The stack normal line is the associated line
\begin{equation}\label{eq:cutnormal}
 \normal_{\cM_d/\cM^{\leq_\epsilon d}}
 =[(M^{\sst}(d)\times\charline{\epsilon})/\Gm].
\end{equation}
Here $\leq_\epsilon$ means $\leq$ for $\epsilon=+1$ and $\geq$ for
$\epsilon=-1$. Formula~\eqref{eq:cutnormal} follows by taking the normal
bundle of $y=0$ before quotienting. It does not require a global
identification of a whole boundary neighbourhood with a normal bundle.

If the boundary stabiliser is $\mu_r$, the two normal characters are
$\zeta$ and $\zeta^{-1}$. They are dual lines. The polarisation is instead
determined by $L$ and the level twist. There is no general reason for it
to equal the normal line. The ordinary coarse normal sheaf is a third
object; the root example already proves that pulling it back need not
give the stack normal line.

\paragraph{A distinction from residual symmetry.}
The stabiliser in \eqref{eq:cutstabs} belongs to the group being
\emph{quotiented}. A separate residual torus can act on the resulting cut.
A fixed point of that residual action is not automatically an additional
automorphism of the object in the stack. Isotropy of the quotient stack
and fixed loci of a residual action are different data.

\section{One complete cutting toy model}
\label{sec:toycut}

Let
\[
 M=\A^1_x,\qquad t\cdot x=t^n x,\qquad L=\OO_M
\]
with trivial geometric fibre action on $L$. This is an affine example:
its affinisation is the identity morphism and is therefore projective.
The trivial bundle is ample on
the affine scheme. The action contracts to $0$. The only fixed component
is $0$, with $\Phi(0)=0$; the level $d=1$ is regular.

If $h$ denotes the constant degree-one section of $L$, then
$R=k[x,h]$ with tensor degrees $\deg x=0,\deg h=1$ and section weights
$\operatorname{wt}x=n,\operatorname{wt}h=0$. In particular, this is an
example where the degree-zero ring is nonconstant. Suppressing $h$ and
retaining the twist symbol, invariant sections of tensor degree $m$ have
the form $x^a y^b\one_1^m$ with
\begin{equation}\label{eq:toyinvariance}
 na+\epsilon b=m,\qquad a,b\ge0.
\end{equation}

\subsection{The lower cut is the weighted projective line}
For $\epsilon=+1$, the invariant sections are nonzero somewhere whenever
$(x,y)\ne(0,0)$: if $y\ne0$ use $y\one_1$; if $x\ne0$ use
$x\one_1^n$. All positive-degree sections vanish at the origin. Hence
\[
 X^{\sst}_+(1)=\A^2\setminus0,
 \qquad\cM^{\le1}=\PP(n,1).
\]
The invariant algebra is a polynomial ring in
\[
 A=x\one_1^n,\qquad B=y\one_1,
 \qquad \deg_m A=n,\quad\deg_m B=1.
\]
Its $n$th Veronese is $k[A,B^n]$; after regrading, both generators have
degree one. Therefore
\[
 M^{\le1}=\Proj k[A,B]\simeq\PP^1.
\]
This computes the stack cut, coarse cut, and grading in the same example.

Where $y\ne0$ the quotient has coordinate $x/y^n$ and is $\A^1$.
The added boundary is $B\mu_n$, with coarse point $[1:0]$. Its normal
character is $+1$. The cut polarisation has fibre character $+1$ as well
in this particular example, so it is $\OO_{\PP(n,1)}(1)$ and does not
descend when $n>1$. Its $n$th power does descend.

\subsection{The upper cut is an affine root stack}
For $\epsilon=-1$, equation~\eqref{eq:toyinvariance} is $na-b=m$.
Positive tensor degree forces $a>0$, so $x=0$ is unstable. Conversely,
if $x\ne0$, the section $x\one_1^n$ is nonzero. Thus
\[
 X^{\sst}_-(1)=\mathbb G_{m,x}\times\A^1_y,
 \qquad
 \cM^{\ge1}=[(\Gm\times\A^1)/\Gm],
\]
where $t\cdot(x,y)=(t^n x,t^{-1}y)$. Setting $x=1$ étale locally
leaves $\mu_n$ acting on the normal coordinate by $y\mapsto\zeta^{-1}y$:
\[
 \cM^{\ge1}\simeq[\A^1_y/\mu_n],\qquad
 M^{\ge1}\simeq\A^1_v,\quad v=xy^n.
\]
For a ring check, the $n$th Veronese of $B^-$ is generated over $k[v]$
by $A=x\one_1^n$ in regraded degree one. Indeed a monomial in degree
$nr$ is $x^{r+j}y^{nj}\one_1^{nr}=A^r v^j$.
Hence its projective spectrum is $\Spec k[v]$.

The interior $y\ne0$ is $\mathbb G_{m,v}$, while the boundary is again $B\mu_n$
at $v=0$. This upper cut is not proper over $k$. Its stack normal line
has character $-1$, the dual of the lower-cut normal line. The same
polarisation twist has character $+1$ on both boundary stacks. Thus on
the upper cut it is the dual of the boundary normal line.

\begin{center}
\begin{tabular}{@{}lll@{}}
\toprule
 & Lower cut & Upper cut\\
\midrule
Semistable locus & $\A^2\setminus0$ & $\Gm\times\A^1$\\
Stack & $\PP(n,1)$ & $[\A^1/\mu_n]$\\
Coarse space & $\PP^1$ & $\A^1$\\
Interior & $\A^1$ & $\Gm$\\
Boundary stack & $B\mu_n$ & $B\mu_n$\\
Normal character & $+1$ & $-1$\\
Boundary polarisation character & $+1$ & $+1$\\
\bottomrule
\end{tabular}
\end{center}
The boundary identifications use the same subgroup $\mu_n\subset\Gm$.
Keeping this identification is what makes the two normal characters
meaningfully opposite.

\subsection{Weights, sections, and the boundary contribution}
On $M$, the monomial $x^a$ has weight $na$. The lower-cut degree-$m$
sections are those with $na\le m$:
\[
 H^0(\cM^{\le1},\cL^{\otimes m})
 \simeq\operatorname{span}_k\{1,x,\ldots,x^{\lfloor m/n\rfloor}\}
 \quad(m\ge0).
\]
Explicitly $x^a$ on the right corresponds to
$x^a y^{m-na}\one_1^m$ on the left. The displayed isomorphism is the
polynomial identification of these vector spaces; it keeps the original
residual weights $na$ if the residual action scales $x$ alone.

On the boundary, $\cL^{\otimes m}$ has character $m\bmod n$, so
\begin{equation}\label{eq:boundarysections}
 H^0(B\mu_n,\cL^{\otimes m}|_{B\mu_n})
 =\begin{cases}k,&n\mid m,\\0,&n\nmid m.\end{cases}
\end{equation}
For $n=2$, this gives the table
\[
\begin{array}{c|rrrrrrr}
 m&0&1&2&3&4&5&6\\\hline
 \dim H^0(\cM^{\le1},\cL^m)&1&1&2&2&3&3&4\\
 \dim H^0(B\mu_2,\cL^m)&1&0&1&0&1&0&1.
\end{array}
\]
For $m\ge1$, multiplication by $y$ embeds the degree-$(m-1)$
polynomials into the degree-$m$ polynomials. The quotient has a monomial
with $y$-exponent zero exactly when $n\mid m$, namely $x^{m/n}$.
This proves the boundary contribution directly, without invoking any
cohomological vanishing theorem. On sheaves the underlying divisor
sequence is
\[
 0\to\OO_{\cP}(m-1)\xrightarrow{\ y\ }
 \OO_{\cP}(m)\to\OO_{\cD}(m)\to0.
\]

Replacing the boundary stack by its coarse point and replacing its
character line by an ordinary one-dimensional fibre would incorrectly
give one boundary section for \emph{every} $m$. The periodic behaviour in
\eqref{eq:boundarysections} is a first elementary trace of the
representation data relevant to orbifold index calculations.

\subsection{What changes at the critical level}
Keep the same $M,L$, take $d=0$, and choose $\epsilon=+1$. Invariance is
$na+b=0$, so the only invariant sections are constant in $x,y$.
Their positive tensor powers are nonzero everywhere. Therefore
\[
 X^{\sst}_+(0)=\A^2,\qquad
 \cM^{\le0}=[\A^2/\Gm],\qquad M^{\le0}=\Spec k.
\]
The origin has stabiliser $\Gm$, so this stack is not DM. It also has
many distinct geometric orbits mapping to the same point. Thus the point
is a good moduli space, not a coarse moduli space of the whole quotient
stack. This calculation isolates the two changes at a wall: continuous
stabilisers can appear, and the GIT quotient can cease to classify
geometric isomorphism classes.

\section{Exercises with solutions}
\label{sec:exercises}

\begin{exercise}\label{ex:trivial}
Compare $[\A^1/\mu_n]$ for the trivial action and for scalar multiplication.
Find their coarse spaces, stabilisers, and the open substacks away from zero.
\end{exercise}
\begin{exercise}\label{ex:torsor}
Show that $[\Gm/\mu_n]\simeq\Gm$ for scalar multiplication, and explain why
this does not extend to an isomorphism $[\A^1/\mu_n]\simeq\A^1$.
\end{exercise}
\begin{exercise}\label{ex:bline}
List the line bundles on $B\mu_3$, compute their global sections, and determine
which descend to the coarse point. What are their tensor cubes?
\end{exercise}
\begin{exercise}\label{ex:weightcount}
On $\PP(3,1)$, list the monomials in degrees $2,3,4,6$. Compute the
pushforwards of $\OO(2),\OO(3),\OO(-1)$. Which of these stack lines descend?
\end{exercise}
\begin{exercise}\label{ex:fibre}
For $[\A^1_z/\mu_n]\to\A^1_u$, calculate the scheme-theoretic fibre over
$u=0$ and its reduction. Calculate the normal and conormal characters of
the reduced origin substack $B\mu_n$ inside the ambient stack
$[\A^1_z/\mu_n]$.
\end{exercise}
\begin{exercise}\label{ex:generic}
Let $\Gm$ act on $\A^2\setminus0$ with weights $(n,n)$. Show that the
quotient is DM but has stabiliser $\mu_n$ at every geometric point. Why
does this differ from $\PP(n,1)$? Does the coarse space detect the difference?
\end{exercise}
\begin{exercise}\label{ex:push}
In the quadratic-cone example, verify explicitly that the odd module $M$
is generated by $x,y$, has generic rank one, and is not locally free at
the cone vertex. Explain why a stack line with this pushforward cannot
be the pullback of a line bundle on the cone.
\end{exercise}
\begin{exercise}\label{ex:inertia}
The \emph{inertia stack} $I\cX$ parameterises pairs $(\xi,\alpha)$ with
$\alpha\in\Aut(\xi)$. Compute $I[\A^1/\mu_n]$ for scalar multiplication.
Which components come from the nonidentity automorphisms?
\end{exercise}
\begin{exercise}\label{ex:twist}
In the lower-cut toy model replace $d=1$ by a positive integer $d$.
Determine the semistable locus, the stack, its polarisation, and the
criterion for that polarisation to descend. What changes when $n=1$?
\end{exercise}
\begin{exercise}\label{ex:twonormals}
For $n=3$, restrict the two stack normal lines and the two polarisation
lines in the cutting toy model to the common boundary $B\mu_3$.
Write their characters. Which have nonzero invariant vectors?
\end{exercise}

\subsection*{Solutions}
\begin{enumerate}[label=\textbf{\arabic*.},leftmargin=*]
\item For Exercise~\ref{ex:trivial}, the trivial-action stack is
$\A^1\times B\mu_n$ with coarse $\A^1$ and stabiliser $\mu_n$
everywhere. Its restriction to $\Gm$ is $\Gm\times B\mu_n$.
For scalar multiplication, the invariant ring is $k[z^n]$; the coarse
space is again $\A^1$, but only the origin has stabiliser $\mu_n$.
The nonzero open is the ordinary scheme $\Gm$.

\item In Exercise~\ref{ex:torsor}, $z\mapsto z^n$ is a $\mu_n$-torsor on
$\Gm$, so the free-quotient calculation gives the isomorphism.
At zero there are $n$ automorphisms in the quotient stack, whereas a
scheme point has only the identity. Equivalence of stacks preserves
automorphism groups, so the proposed extension is impossible.

\item In Exercise~\ref{ex:bline}, the characters are $0,1,2$ modulo $3$.
Their section spaces are $k,0,0$, and only character $0$ descends.
All three tensor cubes are trivial. A torsion line on a stack can thus
be nontrivial even when its coarse space is a point.

\item In Exercise~\ref{ex:weightcount}, the bases are
\[
 \begin{array}{c|l}
 2&y^2\\
 3&y^3,x\\
 4&y^4,xy\\
 6&y^6,xy^3,x^2.
 \end{array}
\]
Formula~\eqref{eq:weightedpush} gives $\OO$, $\OO(1)$, and
$\OO(-1)$ on the coarse $\PP^1$. Of the three stack lines listed,
only $\OO(3)$ descends. In particular, the invertible pushforward of
$\OO(-1)$ is not proof of descent.

\item In Exercise~\ref{ex:fibre}, the fibre is
$[\Spec k[z]/(z^n)/\mu_n]$ and the reduction is $B\mu_n$.
The normal character is $+1$ and the conormal character is $-1$.
The nilpotent thickening records the multiplicity $n$ in
$\pi^*(0)=n\cD$.

\item In Exercise~\ref{ex:generic}, any nonzero coordinate forces
$t^n=1$, so every stabiliser is $\mu_n$. The coarse space is $\PP^1$,
as follows from the invariant ratio on each chart or regrading the
homogeneous coordinate ring. The difference from $\PP(n,1)$ is generic
isotropy. This quotient is a $\mu_n$-gerbe over $\PP^1$; this description
does not assert a product $\PP^1\times B\mu_n$. A global product
would require a trivialisation of the gerbe.

\item For Exercise~\ref{ex:push}, an odd monomial of positive $x$-exponent
is $x$ times an even monomial; if its $x$-exponent is zero, it is $y$
times an even monomial. Thus $M=Ax+Ay$. Generically,
$y=(xy/x^2)x$, so the rank is one. Modulo the vertex ideal the two
degree-one generators remain independent. Hence the module is not
locally free of rank one there. Moreover its fibre character at the
stack origin is sign, which directly obstructs descent. If it were
pulled back from a coarse line, its pushforward would recover that line
and be locally free, giving a second contradiction.

\item In Exercise~\ref{ex:inertia}, a geometric pair is $(z,\zeta)$ with
$\zeta z=z$. Since $\mu_n$ is abelian, quotienting each fixed locus
does not permute the labels $\zeta$. Thus
\[
 I[\A^1/\mu_n]
 \simeq[\A^1/\mu_n]\ \sqcup\!
 \coprod_{\zeta\ne1}B\mu_n.
\]
The identity component is the original stack; each other automorphism
has fixed locus $0$ and supplies a separate $B\mu_n$ component. This
is a first example of the extra sectors used in orbifold formulas.

\item In Exercise~\ref{ex:twist}, invariance is $na+b=md$. For any
$d>0$, both axes away from the origin have some nonvanishing invariant
positive-degree section, while the origin does not. Hence the
semistable locus and stack remain $\A^2\setminus0$ and $\PP(n,1)$.
The polarisation is now $\OO_{\PP(n,1)}(d)$; it descends exactly when
$n\mid d$. The coarse space remains $\PP^1$. When $n=1$ the action
on $\A^2\setminus0$ is free, the stack is $\PP^1$, and every such
polarisation is the ordinary $\OO_{\PP^1}(d)$.

\item In Exercise~\ref{ex:twonormals}, the lower normal has character
$1$, the upper normal has character $-1=2$ modulo $3$, and both
polarisations have character $1$. None of these four boundary lines
has a nonzero invariant vector. Their cubes are all trivial character
lines and each has a one-dimensional invariant space.
\end{enumerate}

\section{A reading plan}
\label{sec:dictionary}

For a second pass, work through the sections in the following order.
\begin{enumerate}
\item Recompute the two affine invariant rings in Section~\ref{sec:quotients}.
Then explain the torsor $s^n=t$ to yourself as a family that cannot be
described by a global root. These are the two independent reasons for
retaining stabilisers and descent.
\item Read the quotient-stack definition alongside the atlas identity
$T\times_{[X/G]}X=P$. This identity explains most of the formal vocabulary.
Use the diagonal-as-isomorphisms description to understand the DM criterion.
\item Work out $\PP(3,1)$ on both charts, including the characters of
$\OO(1)$ and the normal line. Verify the floor in the pushforward formula
and the failure of tensor compatibility.
\item Return to the cut. Keep separate the tensor degree $m$, the group
weight $w$, the stabiliser character modulo its order, and any residual
weight. Derive \eqref{eq:cutinvariance} before using an inequality.
\end{enumerate}

For further reading, Casalaina-Martin--Wise \cite[\S\S2--6]{CMW} develops
the categorical language with many examples. Behrend
\cite[\S\S3.3--3.6]{Behrend} is particularly useful for the quotient,
bundle, and root-stack geometry used here. Vistoli \cite{Vistoli}
provides a more systematic treatment of descent. The Stacks Project
provides precise definitions and theorem statements; its tag references
are more stable than section numbering. Alper \cite{Alper} is the next
step when critical GIT levels and vector-bundle descent become central.
Cadman \cite{Cadman} is the original source for the divisor root construction.

\appendix
\section{Relation to the repository and source conventions}

These notes are a separate introductory document, not proposed manuscript
edits. They use the repository as motivation and as a source of notation.
The mathematical definitions are stated independently, and the examples
are worked out in the text. General existence, local-structure, and
descent theorems have been identified as such rather than claimed to
have elementary proofs here.

The project-specific points of contact are:
\begin{itemize}
\item \path{claude-notes/conventions.tex}, especially its section on
cutting data and names, supplies the convention
$w+\epsilon\ell-md=0$, calligraphic stack notation, and the names of
the two cuts.
\item \path{claude-notes/comparison-theorem.tex}, particularly the
regular/critical-level remarks and the boundary discussion, supplies
the context in which finite boundary stabilisers and polarisation
pushforward matter. The general comparison theorem is not needed for
the elementary toy calculation.
\item \path{SemiprojectiveAlgebraicCutting-integrated.tex} discusses
stack normal lines, finite quotient slices, and possible failure of
coarse polarisation descent. The examples here explain those phenomena
from the beginning.
\end{itemize}

The older integrated manuscript uses a different sign convention for
section weights and the auxiliary direction. The current convention
note records the translation $w_{\mathrm{old}}=-w$ and
$\epsilon_{\mathrm{old}}=-\epsilon$; its old $Y_-$ corresponds to the
current $M^{\le d}$. Thus a cut label or inequality from that older text
must be translated before comparing it with these notes. The invariant
equation in Section~\ref{sec:cutting} is the convention used throughout
this document.

The meaning of ``semiprojective'' also varies in the literature. The
projective-affinisation hypothesis used here is stated separately from
the existence of a contracting action. Nothing in the definition of a
DM stack depends on either property.

\begin{thebibliography}{99}
\bibitem{CMW}
S.~Casalaina-Martin and J.~Wise,
\emph{An introduction to moduli stacks, with a view towards Higgs bundles
on algebraic curves}, arXiv:1708.08124.
\url{https://arxiv.org/abs/1708.08124}.
The saved 147-page PDF in the local Zotero library was consulted;
section and page locators refer to that version.

\bibitem{Behrend}
K.~Behrend, \emph{Introduction to algebraic stacks}, lecture notes,
local Zotero record dated 2014, 127-page version.
\url{https://personal.math.ubc.ca/~behrend/math615A/stacksintro.pdf}.
Relevant material: \S\S3.3--3.6, pp.~110--125. We use the modern coarse
moduli definition given in Section~\ref{sec:coarse}; older conventions
for this terminology differ.

\bibitem{Vistoli}
A.~Vistoli, \emph{Notes on Grothendieck topologies, fibered categories
and descent theory}, arXiv:math/0412512; also in
\emph{Fundamental Algebraic Geometry}, Mathematical Surveys and
Monographs 123, American Mathematical Society, 2005.
\url{https://arxiv.org/abs/math/0412512}.

\bibitem{Stacks}
The Stacks Project Authors, \emph{The Stacks Project}.
\url{https://stacks.math.columbia.edu/}.
Tags used here: 026N (definitions), 04UV (quotient stacks),
04TK (groupoid quotients), 06N3 (DM criterion), and 04UX (coarse-space
literature). Consulted 13 November 2026.

\bibitem{Alper}
J.~Alper, \emph{Good moduli spaces for Artin stacks},
Annales de l'Institut Fourier \textbf{63} (2013), 2349--2402.
\url{https://aif.centre-mersenne.org/articles/10.5802/aif.2833/}.
For vector-bundle descent see Theorem~10.3, journal p.~2386.

\bibitem{Cadman}
C.~Cadman, \emph{Using stacks to impose tangency conditions on curves},
American Journal of Mathematics \textbf{129} (2007), 405--427;
arXiv:math/0312349.
\url{https://arxiv.org/abs/math/0312349}.
For the root construction see \S2.
\end{thebibliography}

\clearpage
\section{Glossary of terms and notation}
\label{sec:glossary}

The entries are alphabetical. They use the same characteristic-zero base
field and conventions as the main text; section references lead to fuller
explanations and worked examples. A notation table follows the definitions. Each entry links to its examples
in Appendix~\ref{sec:examples-glossary}, which in turn links back here.
Closely related terms have separate entries when confusing them would
change a geometric statement.

\newcommand{\glossentry}[3][]{%
  \par\smallskip\phantomsection\label{gloss:#1}%
  \noindent\textbf{#2.}\ #3\ %
  \textit{\hyperref[examples:#1]{Examples (Appendix~\ref*{sec:examples-glossary}).}}\par}

\subsection*{A--C}

\glossentry[affine-invariant-quotient]{Affine invariant quotient}{For an affine scheme $\Spec A$
with a group action, the scheme $\Spec A^G$ defined by invariant functions.
It can identify different orbits, so it need not classify geometric
isomorphism classes of the quotient stack. See
Sections~\ref{sec:quotients} and~\ref{sec:coarse}.}

\glossentry[affinisation]{Affinisation}{The canonical map
$M\to\Spec\Gamma(M,\OO_M)$, which gathers all global regular functions
on $M$ into one morphism to an affine scheme. Projectivity of this map
is the semiprojectivity condition used here; it is not the same as
projectivity of $M$ over $k$. See Section~\ref{sec:cutting}.}

\glossentry[algebraic-space]{Algebraic space}{A sheaf of sets for the étale topology with
diagonal representable by schemes and an étale surjective cover by a
scheme. It generalises a scheme without introducing nontrivial
automorphisms of its objects. Every scheme is an algebraic space, but
not every algebraic space is a scheme. See Section~\ref{sec:dm}.}

\glossentry[algebraic-stack]{Algebraic stack}{A stack in groupoids whose diagonal is
representable by algebraic spaces and which has a representable smooth
surjective atlas from a scheme. The diagonal condition makes
isomorphisms between families algebraic objects; the atlas gives
scheme charts. See Section~\ref{sec:dm}.}

\glossentry[ample-line-bundle]{Ample line bundle}{A line bundle whose positive tensor powers
provide sufficiently many sections to describe the scheme by affine
nonvanishing loci; on a projective scheme, a sufficiently high power
gives a projective embedding. In these notes ampleness is used on the
scheme before taking the GIT quotient. The trivial line bundle on an
affine scheme is ample.}

\glossentry[artin-stack]{Artin stack}{The term commonly used for an algebraic stack
defined using a smooth atlas. It allows positive-dimensional stabilisers,
as in $B\Gm$, and includes DM stacks as a subclass. It does not mean
that every stabiliser must be positive-dimensional. See Section~\ref{sec:dm}.}

\glossentry[atlas]{Atlas}{A representable smooth surjective morphism
$U\to\cX$ from a scheme. An étale atlas is one which is étale, and
exists precisely for DM stacks. Surjectivity permits local lifts of
families; it does not guarantee a global lift of every family to $U$.
See Sections~\ref{sec:quotientstacks} and~\ref{sec:dm}.}

\glossentry[automorphism]{Automorphism}{An isomorphism of an object with itself.
The automorphisms of an object form a group. For a geometric object of
a quotient stack represented by $x\in X$, this is the stabiliser
$\{g\in G:gx=x\}$.}

\glossentry[boundary-stack-and-coarse-boundary]{Boundary stack and coarse boundary}{At a regular cutting
level, the boundary
stack is the closed substack $y=0$ before forgetting stabilisers;
its coarse boundary is the corresponding quotient space. In the running
example they are $B\mu_n$ and a single ordinary point. At critical
levels, the quotient need only be a good moduli space. See
Sections~\ref{sec:cutting} and~\ref{sec:toycut}.}

\glossentry[cartier-divisor]{Cartier divisor}{An effective Cartier divisor is a closed
subscheme locally defined by one non-zero-divisor. It has an associated
line bundle $\OO(D)$ and a canonical section whose zero divisor is $D$.
The same definition is made on a stack by checking it on scheme charts.
See Section~\ref{sec:roots}.}

\glossentry[character]{Character}{A homomorphism of algebraic groups $G\to\Gm$,
equivalently a one-dimensional representation. Characters of $\mu_n$
are $\zeta\mapsto\zeta^c$, indexed by $c$ modulo $n$. Characters of
$\Gm$ are $t\mapsto t^c$, indexed by integers. See Section~\ref{sec:sheaves}.}

\glossentry[classifying-stack]{Classifying stack}{The stack $BG=[\Spec k/G]$ assigning to
$T$ the groupoid of $G$-torsors over $T$. In particular, $B\Gm$
classifies line bundles and $B\mu_n$ classifies pairs
$(N,N^{\otimes n}\simeq\OO_T)$. See Section~\ref{sec:quotientstacks}.}

\glossentry[coarse-moduli-space]{Coarse moduli space}{A map $\pi:\cX\to X$ to an algebraic
space which is universal for maps to algebraic spaces and gives a
bijection between geometric isomorphism classes and geometric points
of $X$, over every algebraically closed extension. It forgets
automorphisms and need not preserve families over general schemes.
See Section~\ref{sec:coarse}.}

\glossentry[cocycle-condition]{Cocycle condition}{The identity
$\alpha_{ij}\alpha_{jk}=\alpha_{ik}$ on triple overlaps for chosen gluing
isomorphisms. For transition functions of a torsor it reads
$g_{ij}g_{jk}=g_{ik}$. It is the compatibility needed to glue local
objects. See Section~\ref{sec:families}.}

\glossentry[conormal-line]{Conormal line}{For a Cartier divisor with ideal sheaf
$\mathcal I$, the line bundle $\mathcal I/\mathcal I^2$ on the divisor.
It is dual to the normal line. Thus if a stabiliser acts on a normal
vector by $\zeta$, it acts on a conormal covector by $\zeta^{-1}$.
See Section~\ref{sec:roots}.}

\glossentry[contracting-action]{Contracting action}{In the sense relevant here, a
$\Gm$-action for which $\lim_{t\to0}t\cdot p$ exists in $M$ for every
geometric point $p$. The limits lie in the fixed locus. Existence of
these limits is an extra hypothesis, distinct from projectivity of
the affinisation. See Section~\ref{sec:cutting}.}

\glossentry[critical-cutting-level]{Critical cutting level}{A level $d$ equal to $\Phi(F)$ for
some fixed component $F$, where $\Phi(F)$ is its section weight.
At such a level fixed points can survive semistability and give
$\Gm$ stabilisers, so the quotient stack can fail to be DM.
See Sections~\ref{sec:cutting} and~\ref{sec:toycut}.}

\subsection*{D--F}

\glossentry[deligne-mumford-stack]{Deligne--Mumford stack}{An algebraic stack admitting a
representable étale surjective atlas from a scheme; equivalently, its
diagonal is unramified. The separated finite-type stacks and finite-type
quotient examples used here have finite reduced stabilisers. With the
broad Stacks Project convention, DM alone does not impose
quasi-compactness of the diagonal or finiteness of every stabiliser.
See Section~\ref{sec:dm}.}

\glossentry[descent-datum]{Descent datum}{Objects on a cover, with specified
isomorphisms on overlaps satisfying a cocycle condition. Such a datum
is effective when it comes from an object on the base. A stack
requires both effectiveness for objects and sheaf-like gluing for
isomorphisms. See Section~\ref{sec:families}.}

\glossentry[descent-of-a-line-bundle-to-the-coarse-space]{Descent of a line bundle to the coarse space}{The existence
of a line bundle $L$ on $X$ and an isomorphism $\cL\simeq\pi^*L$.
In the separated finite-type DM setting here, this holds exactly when
all stabilisers act trivially on the fibres of $\cL$. It is stronger
than invertibility of $\pi_*\cL$. See Section~\ref{sec:sheaves}.}

\glossentry[diagonal]{Diagonal}{The morphism
$\Delta_{\cX}:\cX\to\cX\times\cX$. Its pullback along a pair of
families is their sheaf of isomorphisms. Properties of the diagonal
therefore control isomorphisms, automorphisms, and separation.
See Section~\ref{sec:dm}.}

\glossentry[dimension-of-a-quotient-stack]{Dimension of a quotient stack}{For the equidimensional
examples with a smooth group, $\dim[X/G]=\dim X-\dim G$.
Thus $\PP(n,1)$ has dimension one although its standard atlas is
two-dimensional. The dimensions of stabilisers need not equal the
dimension of the group used in a presentation. See Section~\ref{sec:dm}.}

\glossentry[effective-descent]{Effective descent}{The assertion that compatible local
objects with specified gluing isomorphisms arise from a global object.
For a stack, the global object is unique up to an isomorphism compatible
with the gluing data. This does not remove its automorphisms.
See Section~\ref{sec:families}.}

\glossentry[equivariant-map]{Equivariant map}{A map commuting with the declared group
actions. With left actions, $f(gp)=g f(p)$. Quotient-stack objects
are torsors equipped with such maps; changing to right torsors requires
the corresponding inverse convention. See Section~\ref{sec:quotientstacks}.}

\glossentry[equivariant-vector-bundle]{Equivariant vector bundle}{A vector bundle with a compatible
lift of the group action, linear on the fibres. It is equivalently a
vector bundle on the quotient stack. At a fixed point its fibre is a
representation of the stabiliser. In rank one the lift is a
linearisation and its fibre representation is a character.
See Section~\ref{sec:sheaves}.}

\glossentry[etale-cover]{Étale cover}{A jointly surjective family of étale morphisms.
Its overlaps are fibre products over the base. Such a cover can permit
choices, such as roots of units, which are unavailable on any Zariski
open cover. See Section~\ref{sec:families}.}

\glossentry[etale-morphism]{Étale morphism}{A flat, unramified morphism locally of
finite presentation. It is the algebraic model for a local isomorphism
without branching. For example, adjoining an $n$th root of a unit
is finite étale when $n$ is invertible. See Section~\ref{sec:families}.}

\glossentry[family]{Family}{An object parameterised by a test scheme $T$;
in stack language it is an object of $\cX(T)$, or a morphism
$T\to\cX$. Its pullback along $T'\to T$ is the family with changed
base. Testing only over $\Spec k$ loses information about families.
See Section~\ref{sec:families}.}

\glossentry[fibre-product-of-stacks]{Fibre product of stacks}{A 2-fibre product: its objects are
two objects together with a specified isomorphism between their images
in the common target. This extra isomorphism is why
$X\times_{[X/G]}X\simeq G\times X$. See
Section~\ref{sec:quotientstacks}.}

\glossentry[flat-and-faithfully-flat]{Flat and faithfully flat}{A morphism of schemes is flat if
each local-ring map makes the target local ring a flat module: tensoring
with that module preserves exact sequences. A flat morphism which is
surjective is faithfully flat. Flatness is part of the definitions of
étale and fppf covers.}

\glossentry[fppf-topology]{Fppf topology}{The topology whose covers are jointly
surjective families of flat morphisms locally of finite presentation.
The abbreviation comes from the French for faithfully flat and of
finite presentation. It contains the étale topology. The stacks in
these notes satisfy descent in this topology. See
Section~\ref{sec:families}.}

\glossentry[free-action]{Free action}{An action with trivial stabilisers; for the
finite-type smooth-group actions over characteristic zero considered
here, this can be checked on geometric points. If $X\to Y$ is a
$G$-torsor, then $[X/G]\simeq Y$. Freeness is stronger than the finite
stabiliser condition used in GIT stability. See
Section~\ref{sec:quotientstacks}.}

\glossentry[functor-of-points]{Functor of points}{The functor
$T\mapsto\Hom(T,X)$, together with pullback along maps of test schemes.
It determines a scheme $X$. For a stack, the analogous construction
has groupoids as values, so it retains isomorphisms as well as objects.
See Section~\ref{sec:families}.}

\subsection*{G--L}

\glossentry[geometric-fibre-weight]{Geometric fibre weight}{The integer $c$ describing an
actual action on a line fibre by $v\mapsto t^c v$. For $\mu_n$ only
its residue modulo $n$ matters. Under the section-action convention
here, the fibre weight of $L$ at a fixed component is $-\Phi(F)$,
while the nonvanishing section weight is $\Phi(F)$.
See Sections~\ref{sec:sheaves} and~\ref{sec:cutting}.}

\glossentry[geometric-point]{Geometric point}{A morphism $\Spec\Omega\to X$ where
$\Omega$ is an algebraically closed field. For a stack this is an
object over $\Omega$ and can have nontrivial automorphisms. Passing
to its geometric isomorphism class is the step used in defining a
coarse moduli space. See Section~\ref{sec:coarse}.}

\glossentry[git-quotient]{GIT quotient}{The quotient determined by a chosen
linearisation and its invariant sections. In the projective or relative
construction here it is $\Proj B$ in tensor degree. At a regular
level it is the coarse space of the semistable quotient stack; at a
critical level it can identify distinct geometric orbit classes.
See Sections~\ref{sec:coarse} and~\ref{sec:cutting}.}

\glossentry[good-moduli-space]{Good moduli space}{A quasi-compact map
$\pi:\cX\to X$ to an algebraic space such that pushforward on
quasi-coherent sheaves is exact and
$\OO_X\simeq\pi_*\OO_{\cX}$. It can identify distinct geometric
isomorphism classes, unlike a coarse moduli space. The map
$[\A^1/\Gm]\to\Spec k$ for scaling is an example.
See Section~\ref{sec:coarse}.}

\glossentry[groupoid]{Groupoid}{A category in which every arrow is invertible.
The action groupoid has objects $x\in X$ and arrows $g:x\to gx$.
Its isomorphism classes are orbits, while its automorphism groups are
stabilisers. See Section~\ref{sec:quotients}.}

\glossentry[groupoid-presentation]{Groupoid presentation}{An atlas $U\to\cX$ together with
the arrow space $R=U\times_{\cX}U$, source and target maps, and
composition. The stack is recovered by allowing this groupoid to
describe families locally and then glue. For $[X/G]$, the arrow space
is $G\times X$. See Section~\ref{sec:quotientstacks}.}

\glossentry[inertia-stack]{Inertia stack}{The stack $I\cX$ of pairs $(\xi,\alpha)$,
where $\xi$ is an object and $\alpha$ is an automorphism of it.
It records stabilisers in families. For the scalar cyclic quotient,
$I[\A^1/\mu_n]$ has the identity component $[\A^1/\mu_n]$
and one extra $B\mu_n$ component for each nonidentity root of unity.
See Exercise~\ref{ex:inertia}.}

\glossentry[invariant-function-or-section]{Invariant function or section}{A function unchanged by
the action, or a section fixed by the action induced by a
linearisation. For a line with geometric fibre character $\chi$,
equivariance of a section means $f(gx)=\chi(g) f(x)$; invariance does
not require the coefficient function itself to have weight zero.
See Sections~\ref{sec:sheaves} and~\ref{sec:cutting}.}

\glossentry[isomorphism-sheaf]{Isomorphism sheaf}{For two families $\xi,\eta$ over $T$,
the sheaf assigning to $T'\to T$ the isomorphisms between their
pullbacks to $T'$. It is denoted
$\underline{\Isom}_T(\xi,\eta)$ and is a pullback of the diagonal.
When the two families coincide it is their automorphism group sheaf.
See Section~\ref{sec:dm}.}

\glossentry[line-bundle]{Line bundle}{A locally free sheaf of rank one, or the
corresponding geometric bundle with one-dimensional vector-space
fibres. A line bundle on a stack also records how automorphisms act
on those fibres. These fibre characters can obstruct descent to a
coarse space. See Section~\ref{sec:sheaves}.}

\glossentry[linearisation]{Linearisation}{A lift of a group action to a line bundle,
linear on its fibres and compatible with group multiplication. It
specifies an action on sections and hence the invariant sections used
to form a GIT quotient. Different linearisations on the same
underlying bundle can give different stability conditions.
See Section~\ref{sec:cutting}.}

\glossentry[linearly-reductive-group]{Linearly reductive group}{A group for which taking
invariant vectors is an exact operation on representations. Finite
groups in characteristic zero have this property by averaging over
the group. This exactness explains the good behaviour of invariant
modules and tame coarse pushforward here. See Section~\ref{sec:sheaves}.}

\glossentry[locally-free-sheaf]{Locally free sheaf}{A sheaf which is isomorphic locally
to a finite direct sum of copies of the structure sheaf. Rank-one
locally free sheaves are line bundles. A coherent sheaf of generic
rank one need not be locally free at singular points, as the odd
module on the quadratic cone demonstrates. See Section~\ref{sec:sheaves}.}

\subsection*{N--R}

\glossentry[normal-line]{Normal line}{For a Cartier divisor
$\cD\subset\cX$, the line bundle
$\normal_{\cD/\cX}=(\mathcal I/\mathcal I^2)^\vee
\simeq\OO_{\cX}(\cD)|_{\cD}$. It includes the action of every
stabiliser on the normal direction. It need not be the pullback of
the normal line of the coarse divisor. See Section~\ref{sec:roots}.}

\glossentry[orbit]{Orbit}{The points obtained from $x$ by acting with the
group. In a quotient stack an orbit gives one geometric isomorphism
class, with automorphism group $\Stab(x)$. An orbit can fail to be
closed, which matters when comparing a GIT quotient with a coarse
moduli space. See Sections~\ref{sec:quotients} and~\ref{sec:coarse}.}

\glossentry[polarisation]{Polarisation}{The chosen ample linearised line bundle
used in the quotient construction, or the associated line on the
quotient stack. Its pushforward to the coarse space can fail to be
invertible, and even an invertible pushforward need not be a
descended polarisation. It is also distinct from the boundary normal
line. See Sections~\ref{sec:sheaves} and~\ref{sec:cutting}.}

\glossentry[projective-degree]{Projective degree}{The tensor degree $m$ in a section
ring $R=\bigoplus_{m\ge0}H^0(M,L^{\otimes m})$. The operation
$\Proj$ in these notes uses this grading, not the group-weight
grading. Degree-zero functions need not be constant.
See Section~\ref{sec:cutting}.}

\glossentry[pullback]{Pullback}{The operation taking a sheaf or family to a new
base along a morphism. For a line bundle $L$ on a coarse space,
$\pi^*L$ is a line on the stack with trivial stabiliser actions.
Pullback is the operation used in the definition of line-bundle
descent. See Section~\ref{sec:sheaves}.}

\glossentry[pushforward]{Pushforward}{The sheaf whose sections downstairs are
sections upstairs over the inverse image. On affine finite-group
quotient charts it is computed by taking invariant modules.
It can lose stabiliser characters, fail to preserve local freeness,
and fail to commute with tensor products.
See Sections~\ref{sec:sheaves} and~\ref{sec:weighted}.}

\glossentry[quasi-coherent-sheaf]{Quasi-coherent sheaf}{A sheaf of modules which, on affine
charts $\Spec A$, is associated to an $A$-module, with compatible
restriction maps. On a stack this condition is checked on an atlas
with descent data. Vector bundles give quasi-coherent sheaves;
the category also includes invariant modules which are not locally
free. See Section~\ref{sec:sheaves}.}

\glossentry[quotient-stack]{Quotient stack}{The stack $[X/G]$ whose objects over $T$
are $G$-torsors $P\to T$ with equivariant maps $P\to X$.
Arrows are torsor isomorphisms compatible with those maps. The torsor
retains the transition functions needed to glue a family, rather
than requiring one global map $T\to X$. See
Section~\ref{sec:quotientstacks}.}

\glossentry[regular-cutting-level]{Regular cutting level}{A level $d$ unequal to
$\Phi(F)$ for every fixed component $F$. Under the hypotheses
of the cutting discussion, the semistable locus then has closed
orbits within it and finite stabilisers, and the quotient stack is
smooth DM. Regularity does not imply freeness.
See Section~\ref{sec:cutting}.}

\glossentry[representable-morphism]{Representable morphism}{A morphism
$\cX\to\cY$ representable by algebraic spaces is one whose base
change by every scheme $T\to\cY$ is an algebraic space. If every
such base change is a scheme, it is representable by schemes.
Properties such as étaleness of a representable map are tested on
these base changes. See Section~\ref{sec:dm}.}

\glossentry[residual-action]{Residual action}{An additional action that commutes with
the group being quotiented and therefore induces an action on the
quotient. Its fixed points describe an external symmetry of the
resulting stack; they are different data from automorphisms of
objects in that stack. See Section~\ref{sec:cutting}.}

\glossentry[residual-gerbe-at-a-closed-point]{Residual gerbe at a closed point}{The substack retaining
the point and its automorphisms. For a closed point with stabiliser
$H$ over the algebraically closed field here it is $BH$.
In $[\A^1/\mu_n]$ the reduced origin substack is $B\mu_n$,
whereas the scheme-theoretic fibre of the coarse map is thicker.
See Section~\ref{sec:roots}.}

\glossentry[root-divisor]{Root divisor}{The Cartier substack $\cD$ where the
universal section of a divisor root stack vanishes. It satisfies
$\pi^*D=n\cD$. For smooth $D$, it is the reduced substack over
$D$ and is locally $D\times B\mu_n$, although a global product
need not exist. See Section~\ref{sec:roots}.}

\glossentry[root-gerbe]{Root gerbe}{The stack of $n$th roots of a line bundle
$L$, with no section specified: its objects have an isomorphism
$N^{\otimes n}\simeq f^*L$. It has $\mu_n$ stabiliser everywhere.
For $L=\OO_X$ it is $X\times B\mu_n$.
See Section~\ref{sec:roots}.}

\glossentry[root-stack-of-a-divisor]{Root stack of a divisor}{The stack
$\sqrt[n]{(X,D)}$ of data $(f:T\to X,N,s,\phi)$ with
$\phi:N^{\otimes n}\simeq f^*\OO_X(D)$ carrying $s^n$
to the canonical divisor section. It is $X$ away from $D$
and has $\mu_n$ stabilisers over $D$. Its local chart is
$[\Spec A[z]/(z^n-a)/\mu_n]$ if $D=(a)$.
See Section~\ref{sec:roots}.}

\subsection*{S--Z}

\glossentry[scheme-theoretic-fibre]{Scheme-theoretic fibre}{A fibre defined by base change,
including its equations and nilpotents. For the coarse map
$[\A^1_z/\mu_n]\to\A^1_u$ with $u=z^n$, the fibre over
zero is $[\Spec k[z]/(z^n)/\mu_n]$. Its reduction is
$B\mu_n$, but the two are not the same stack.
See Section~\ref{sec:roots}.}

\glossentry[section-ring]{Section ring}{The graded algebra
$R=\bigoplus_{m\ge0}H^0(M,L^{\otimes m})$, with multiplication
given by tensoring sections. A linearisation gives a second,
independent decomposition by group weights $R_{m,w}$.
The finite-generation and $M\simeq\Proj R$ hypotheses are stated
explicitly in the cutting discussion. See Section~\ref{sec:cutting}.}

\glossentry[section-weight]{Section weight}{The integer $w$ in $t\cdot s=t^w s$.
These notes use $(t\cdot s)(p)=t_L^{-1}s(tp)$, so at a fixed
component the weight of a nonvanishing section is the negative
of its geometric fibre weight. Invariant cut sections satisfy
$w+\epsilon\ell-md=0$. See Section~\ref{sec:cutting}.}

\glossentry[semiprojective]{Semiprojective}{In the projective-affinisation sense
used here, a variety whose global-function algebra is finitely
generated and whose affinisation morphism is projective.
Some authors impose additional conditions on a $\Gm$-action,
including existence of contracting limits and properness of the
fixed locus. Those conditions are not silently included here.
See Section~\ref{sec:cutting} and Appendix~A.}

\glossentry[semistable-point]{Semistable point}{For the chosen ample linearisation,
a point where some invariant section of positive tensor degree
does not vanish. The positive-degree requirement matters when
$R_0$ is nonconstant. Semistability depends on the linearisation
and does not by itself imply a finite stabiliser.
See Section~\ref{sec:cutting}.}

\glossentry[separated-stack]{Separated stack}{A stack whose diagonal is proper.
For a DM stack this makes its unramified diagonal finite.
It is a further condition, not part of the broad DM definition
used here. See Section~\ref{sec:dm}.}

\glossentry[smooth-stack]{Smooth stack}{An algebraic stack smooth over the base
field. Smoothness can be checked on a smooth atlas.
Thus $[\A^2/\mu_2]$ for the sign action is smooth even though
its coarse quadratic cone is singular. A DM stack need not
be smooth. See Section~\ref{sec:dm}.}

\glossentry[stabiliser-or-isotropy-group]{Stabiliser or isotropy group}{The automorphism group
of a geometric object. In $[X/G]$, an object represented by
$x$ has stabiliser $G_x=\{g:gx=x\}$. Its representations on
tangent spaces and line fibres are part of stack geometry.
See Sections~\ref{sec:quotients} and~\ref{sec:sheaves}.}

\glossentry[stable-point]{Stable point}{In the standard GIT setting used here,
a semistable point whose orbit is closed within the semistable
locus and whose stabiliser is finite. Stability is weaker than
freeness. Finite stabiliser alone does not ensure that the orbit
is closed. See Section~\ref{sec:cutting}.}

\glossentry[stack]{Stack}{A groupoid-valued moduli problem with pullback,
in which isomorphisms form sheaves and descent data for objects
are effective. Objects glue up to specified compatible
isomorphisms, without discarding their automorphisms.
See Section~\ref{sec:families}.}

\glossentry[tame-stack]{Tame stack}{In the finite-inertia setting, meaning
that $I\cX\to\cX$ is a finite morphism, a stack
whose coarse-space pushforward on quasi-coherent sheaves is exact.
Finite stabilisers are linearly reductive in the characteristic-zero
examples here, so the separated finite-type DM stacks used for
descent are tame. This explains why invariants are exact; it does
not imply descent of every line bundle. See Section~\ref{sec:sheaves}.}

\glossentry[tensor-degree]{Tensor degree}{Another name for the projective degree
$m$, counting powers $L^{\otimes m}$. It is independent of the
group weight $w$ and of any degeneration parameter.
See Section~\ref{sec:cutting}.}

\glossentry[topology-in-descent-theory]{Topology, in descent theory}{A specification of which
families of maps count as covers and hence where objects and
arrows must glue. Zariski, étale, and fppf covers give the
three topologies used here. This is additional cover data on
test schemes, rather than just the usual topology on one
underlying space. See Section~\ref{sec:families}.}

\glossentry[torsor]{Torsor}{A family locally isomorphic to $G\times T$
with the group acting by translation, with no preferred identity
section. A section trivialises the torsor. In our left-action
convention the map $G\times P\to P\times_T P$, $(g,p)\mapsto(gp,p)$,
is an isomorphism. See Section~\ref{sec:quotientstacks}.}

\glossentry[torus]{Torus}{An algebraic group isomorphic over the
algebraically closed field to $(\Gm)^r$. The cutting group is
the one-dimensional torus $\Gm$. Its finite stabilisers are
cyclic groups $\mu_n$; these finite groups are not themselves
positive-dimensional tori. See Section~\ref{sec:cutting}.}

\glossentry[transporter]{Transporter}{For two points or families $x,x'$ in an
action presentation, the space $\{g\in G:gx=x'\}$ of arrows
between them. When $x=x'$ it is the stabiliser. It is a concrete
description of a pullback of the quotient-stack diagonal.
See Section~\ref{sec:dm}.}

\glossentry[unramified-morphism]{Unramified morphism}{For schemes, a morphism locally
of finite type whose relative differential sheaf $\Omega_{X/Y}$
vanishes; geometrically, it has no infinitesimal motion in its
fibres. For algebraic spaces the condition is checked étale
locally. An unramified diagonal characterises DM stacks.
See Section~\ref{sec:dm}.}

\glossentry[vector-bundle]{Vector bundle}{A locally free sheaf of finite rank,
or its associated family of vector spaces. On $[X/G]$, it is
a $G$-equivariant vector bundle on $X$. Its stabiliser
representations must be trivial if it is to descend to the
coarse space in the setting here. See Section~\ref{sec:sheaves}.}

\glossentry[veronese-subring]{Veronese subring}{For a nonnegatively graded ring $B$,
the ring $B^{(N)}=\bigoplus_{r\ge0}B_{Nr}$, regraded so that
$B_{Nr}$ has degree $r$. Its projective spectrum agrees with
$\Proj B$. In the toy lower cut, taking the $n$th Veronese
turns $k[A,B]$ of degrees $(n,1)$ into $k[A,B^n]$
with both generators of regraded degree one.
See Section~\ref{sec:toycut}.}

\glossentry[weight]{Weight}{An integer indexing a character of $\Gm$,
or a character's residue modulo $n$ for $\mu_n$.
One must specify the representation: a section weight,
geometric fibre weight, tangent weight, and residual weight
need not be the same. See Sections~\ref{sec:sheaves}
and~\ref{sec:cutting}.}

\glossentry[weighted-projective-stack]{Weighted projective stack}{The quotient
$[(\A^{r+1}\setminus0)/\Gm]$ for a positive weighted action.
In these notes $\PP(n,1)$ means the stack with action
$(x,y)\mapsto(t^n x,ty)$, whose coarse space is $\PP^1$.
It has a $\mu_n$ stabiliser at $y=0$ and trivial stabilisers
elsewhere. See Section~\ref{sec:weighted}.}

\glossentry[weight-space]{Weight space}{The subspace on which the group acts
by one fixed character. For example, $R_{m,w}$ consists of
sections of $L^{\otimes m}$ with section weight $w$.
Taking invariants selects the weight-zero part of the
declared representation. See Section~\ref{sec:cutting}.}

\glossentry[zariski-versus-etale-local]{Zariski versus étale local}{Zariski local means after
restricting to open subsets. Étale local allows étale covers,
including maps adjoining roots of units. A choice can exist
étale locally without existing on any nonempty Zariski open,
as in the torsor $s^n=t$ on $\Gm$.
See Section~\ref{sec:families}.}

\subsection*{Notation at a glance}

\begin{longtable}{@{}p{0.29\textwidth}p{0.65\textwidth}@{}}
\toprule
Symbol & Meaning\\
\midrule
\endhead
$k$ & An algebraically closed field of characteristic zero.\\
$\Gm$, $\mu_n$ & The multiplicative group and its subgroup of $n$th roots of unity.\\
$\cX$, $X$ & A stack and, when specified, its coarse space.\\
$[X/G]$, $BG$ & A quotient stack and the classifying stack $[\Spec k/G]$.\\
$\pi$ & Usually the map from a stack to its coarse or good moduli space.\\
$\pi^*$, $\pi_*$ & Pullback and pushforward of sheaves.\\
$G_x$, $I\cX$ & The stabiliser at a geometric point and the inertia stack.\\
$\cL$, $L$ & A line bundle on a stack or scheme, respectively, as specified.\\
$\cD$, $D$ & A stack divisor and its coarse divisor.\\
$\normal_{\cD/\cX}$ & The stack normal line of a Cartier divisor.\\
$\sqrt[n]{(X,D)}$ & The $n$th root stack of the divisor pair.\\
$\PP(n,1)$ & The weighted projective stack for point weights $(n,1)$.\\
$\OO_{\PP(n,1)}(c)$ & The line defined by geometric fibre character $t^c$.\\
$m$, $w$ & Tensor degree and original section weight, kept independent.\\
$R_{m,w}$ & The weight-$w$ part of $H^0(M,L^{\otimes m})$.\\
$\Phi(F)$ & The nonvanishing section weight on a fixed component; the geometric fibre weight is $-\Phi(F)$.\\
$\epsilon$, $d$ & Auxiliary point weight $\epsilon=\pm1$ and integral cutting level $d$.\\
$y$, $\one_d$ & Auxiliary coordinate of section weight $\epsilon$ and twist section of weight $-d$.\\
$B^\epsilon$, $D$ & In the cutting algebra, the invariant section ring and threshold ring $\bigoplus_mR_{m,md}$. This use of $D$ is separate from divisor notation.\\
$\cM^{\le d}$, $\cM^{\ge d}$ & The lower and upper cut stacks; roman $M$ denotes their quotients.\\
$\cM_d$, $M_d$ & The common boundary stack and its quotient.\\
\bottomrule
\end{longtable}

For sign calculations, the defining invariant equation is always
\[
 w+\epsilon\ell-md=0,\qquad \ell\ge0.
\]
For descent to a coarse space, the decisive datum is the stabiliser
character on a line fibre, not just the appearance of its pushforward.

\clearpage
\section{Examples Glossary}
\label{sec:examples-glossary}

This appendix gives two small examples for each definition in
Appendix~\ref{sec:glossary}, in the same alphabetical order. Each heading
links to its definition, and each definition links back to its examples.
When one item is a contrasting case, the text states which condition
fails. The examples repeatedly use the same basic spaces so that the
terminology can be learned without a succession of new constructions.

As before, $k$ is algebraically closed of characteristic zero. The
notation $[X/G]$ denotes a stack, and $\PP(2,1)$ denotes the weighted
projective stack, whose coarse space is $\PP^1$. On coordinates,
$t\cdot z=t^a z$ has direct-pullback section weight $a$; a geometric
line fibre acted on by $t^c$ has constant section weight $-c$.

\newcommand{\examplegloss}[3]{%
  \par\medskip\Needspace{5\baselineskip}%
  \phantomsection\label{examples:#1}%
  \noindent\textbf{#2.}\quad
  \hyperref[gloss:#1]{\textit{Glossary entry (Appendix~\ref*{sec:glossary}).}}\par
  #3\par}

\subsection*{A--C}

\examplegloss{affine-invariant-quotient}{Affine invariant quotient}{
\begin{enumerate}[label=(\arabic*),leftmargin=*,itemsep=2pt,topsep=3pt]
\item Let $\mu_2$ act on $\A^1_z$ by $z\mapsto-z$. The invariant polynomials are the even ones, so the quotient is $\Spec k[z^2]\simeq\A^1_u$, with $u=z^2$. The points $z$ and $-z$ have the same image.
\item Let $\Gm$ scale $\A^1_z$. Only constant polynomials are invariant, so the affine invariant quotient is a point. It identifies the zero orbit and the nonzero orbit, although those are distinct geometric isomorphism classes.
\end{enumerate}}

\examplegloss{affinisation}{Affinisation}{
\begin{enumerate}[label=(\arabic*),leftmargin=*,itemsep=2pt,topsep=3pt]
\item For $M=\A^1$, the global-function ring is $k[x]$. Its affinisation $\A^1\to\Spec k[x]$ is the identity. This morphism is projective even though $\A^1$ is not projective over $k$.
\item Every global regular function on $\PP^1$ is constant. Its affinisation is therefore $\PP^1\to\Spec k$, the map to a point. This morphism is projective because $\PP^1$ itself is projective over $k$.
\end{enumerate}}

\examplegloss{algebraic-space}{Algebraic space}{
\begin{enumerate}[label=(\arabic*),leftmargin=*,itemsep=2pt,topsep=3pt]
\item The affine line is an algebraic space: its identity map is an étale atlas, and its diagonal in $\A^1\times\A^1$ is the ordinary closed subscheme $x=x'$. Its points have no nonidentity arrows.
\item The projective line is also an algebraic space. Its two standard affine charts give an étale surjective atlas $\A^1\sqcup\A^1\to\PP^1$. Both examples happen to be schemes; the definition of algebraic space also allows more general étale gluing.
\end{enumerate}}

\examplegloss{algebraic-stack}{Algebraic stack}{
\begin{enumerate}[label=(\arabic*),leftmargin=*,itemsep=2pt,topsep=3pt]
\item The quotient $[\A^1/\mu_2]$ has atlas $\A^1\to[\A^1/\mu_2]$. The atlas is étale, and its arrow scheme is $\mu_2\times\A^1$. The fixed origin retains two automorphisms.
\item The stack $B\Gm$ has smooth atlas $\Spec k\to B\Gm$. Base-changing this atlas by a line bundle $N$ on $T$ gives the torsor $(N^\vee)^\times$ of nonzero dual frames, with our chosen convention. Thus this atlas has one-dimensional fibres, although its source is a point.
\end{enumerate}}

\examplegloss{ample-line-bundle}{Ample line bundle}{
\begin{enumerate}[label=(\arabic*),leftmargin=*,itemsep=2pt,topsep=3pt]
\item On $\PP^1$, the bundle $\OO(1)$ is ample: its coordinate sections $X_0,X_1$ give the usual projective embedding. The nonvanishing loci of these sections are the two affine charts.
\item On $\A^1$, the trivial bundle $\OO$ is ample. Its sections are polynomials, and their nonvanishing loci are principal affine opens. This differs from $\OO$ on $\PP^1$, whose only global sections are constants and which is not ample there.
\end{enumerate}}

\examplegloss{artin-stack}{Artin stack}{
\begin{enumerate}[label=(\arabic*),leftmargin=*,itemsep=2pt,topsep=3pt]
\item The stack $B\Gm$ is Artin: the point is a smooth atlas, and its geometric object has automorphism group $\Gm$. That positive-dimensional group prevents an étale atlas.
\item The scaling quotient $[\A^1/\Gm]$ is Artin. At zero its stabiliser is $\Gm$, but on the nonzero open the action is free and the quotient stack is a point. One Artin stack can therefore contain both continuous-stabiliser and trivial-stabiliser regions.
\end{enumerate}}

\examplegloss{atlas}{Atlas}{
\begin{enumerate}[label=(\arabic*),leftmargin=*,itemsep=2pt,topsep=3pt]
\item The map $\A^1_z\to[\A^1_z/\mu_2]$ is an étale atlas, including at $z=0$. Its base change by a family is the corresponding $\mu_2$-torsor. Its composite to the coarse line, $z\mapsto z^2$, is instead ramified at zero.
\item For $\PP(2,1)$, the map $\A^2\setminus0\to\PP(2,1)$ is a smooth atlas of relative dimension one. The stack is DM, but this particular atlas is not étale; the scheme covers of its smaller quotient charts provide an étale atlas.
\end{enumerate}}

\examplegloss{automorphism}{Automorphism}{
\begin{enumerate}[label=(\arabic*),leftmargin=*,itemsep=2pt,topsep=3pt]
\item A one-dimensional vector space has automorphisms $v\mapsto\lambda v$ for every $\lambda\in k^*$. This explains the automorphism group $\Gm$ of the geometric object of $B\Gm$.
\item In $[\A^1/\mu_2]$, the origin has two automorphisms, represented by $1$ and $-1$. A nonzero point has only the identity automorphism, because $-z=z$ would force $z=0$ in characteristic zero.
\end{enumerate}}

\examplegloss{boundary-stack-and-coarse-boundary}{Boundary stack and coarse boundary}{
\begin{enumerate}[label=(\arabic*),leftmargin=*,itemsep=2pt,topsep=3pt]
\item For the weight-one affine-line cut at level $1$, the lower cut is the ordinary $\PP^1$. The boundary $y=0$ is the point $[1:0]$, with no nontrivial stabiliser, so its stack and coarse descriptions agree.
\item For the weight-two cut, the lower stack is $\PP(2,1)$. Its boundary is $B\mu_2$, while the coarse boundary is still one point of $\PP^1$. The extra sign automorphism acts on the stack normal line.
\end{enumerate}}

\examplegloss{cartier-divisor}{Cartier divisor}{
\begin{enumerate}[label=(\arabic*),leftmargin=*,itemsep=2pt,topsep=3pt]
\item The origin in $\A^1_z$ is an effective Cartier divisor, since $z$ is a non-zero-divisor in $k[z]$. Its ideal is $(z)$ and its divisor line is generated locally by $z^{-1}$.
\item The substack $\{z=0\}\subset[\A^1_z/\mu_2]$ is also Cartier: its pullback to the atlas is the same divisor $(z)$. Its ideal line has an equivariant structure, so taking the quotient does not discard the stabiliser character.
\end{enumerate}}

\examplegloss{character}{Character}{
\begin{enumerate}[label=(\arabic*),leftmargin=*,itemsep=2pt,topsep=3pt]
\item The formula $t\mapsto t^2$ is a character of $\Gm$. A line with this character has geometric action $v\mapsto t^2v$. Tensoring two such lines gives character $t^4$.
\item On $\mu_3$, the characters $\zeta\mapsto\zeta$ and $\zeta\mapsto\zeta^4$ are equal because $\zeta^3=1$. Their common tensor cube is the trivial character. Thus finite stabiliser characters are indexed modulo the group order.
\end{enumerate}}

\examplegloss{classifying-stack}{Classifying stack}{
\begin{enumerate}[label=(\arabic*),leftmargin=*,itemsep=2pt,topsep=3pt]
\item The stack $B\mu_2$ has one geometric isomorphism class with automorphism group $\mu_2$. Over $T=\mathbb G_{m,t}$, it also has the nontrivial family given by the torsor $s^2=t$.
\item A map $\PP^1\to B\Gm$ can specify the line bundle $\OO_{\PP^1}(1)$. The trivial bundle specifies a different map, although the fibres of both bundles over each geometric point are one-dimensional vector spaces. The classifying stack remembers the global gluing.
\end{enumerate}}

\examplegloss{coarse-moduli-space}{Coarse moduli space}{
\begin{enumerate}[label=(\arabic*),leftmargin=*,itemsep=2pt,topsep=3pt]
\item The coarse space of $B\mu_2$ is $\Spec k$: there is one geometric isomorphism class. The coarse point forgets the two automorphisms and the different torsor families over parameter schemes.
\item The coarse map $[\A^1_z/\mu_2]\to\A^1_u$ is $u=z^2$. Each nonzero $u$ corresponds to the single orbit of its two square roots, while $u=0$ corresponds to the origin orbit. Both have one coarse point, despite their different stabilisers.
\end{enumerate}}

\examplegloss{cocycle-condition}{Cocycle condition}{
\begin{enumerate}[label=(\arabic*),leftmargin=*,itemsep=2pt,topsep=3pt]
\item On any three-open cover, choose identity identifications between copies of the trivial line. Then $\alpha_{ij}\alpha_{jk}=\id\,\id=\id=\alpha_{ik}$, so the local bundles glue to the trivial bundle.
\item On a triple overlap, transition functions $g_{12}=2$ and $g_{23}=3$ require $g_{13}=6$. Choosing $g_{13}=5$ would make transporting a vector directly from chart $3$ to chart $1$ disagree with transporting it through chart $2$.
\end{enumerate}}

\examplegloss{conormal-line}{Conormal line}{
\begin{enumerate}[label=(\arabic*),leftmargin=*,itemsep=2pt,topsep=3pt]
\item At the origin of $\A^1_z$, the conormal is $(z)/(z^2)$, a one-dimensional vector space spanned by the class of $z$. Its dual is the usual tangent, or normal, line at that point.
\item For the origin substack in $[\A^1_z/\mu_3]$, a normal vector has geometric character $\zeta$. The conormal covector has the inverse geometric character $\zeta^{-1}=\zeta^2$. Their tensor product has trivial character, as expected for a line and its dual.
\end{enumerate}}

\examplegloss{contracting-action}{Contracting action}{
\begin{enumerate}[label=(\arabic*),leftmargin=*,itemsep=2pt,topsep=3pt]
\item For $t\cdot x=tx$ on $\A^1$, every point has limit $0$ as $t\to0$. Algebraically the formula extends to $(t,x)\mapsto tx$ on $\A^1_t\times\A^1_x$.
\item For $t\cdot(x,y)=(tx,t^2y)$ on $\A^2$, every point again limits to the origin. Both coordinates have positive weights, so both vanish in the limit. The fixed locus is the origin in each example.
\end{enumerate}}

\examplegloss{critical-cutting-level}{Critical cutting level}{
\begin{enumerate}[label=(\arabic*),leftmargin=*,itemsep=2pt,topsep=3pt]
\item On $M=\A^1$ with $t\cdot x=t^2x$ and the trivial fibre action on $L=\OO$, the fixed origin has $\Phi(0)=0$. Thus $d=0$ is critical. In the lower cut every point is semistable, and the product origin has stabiliser $\Gm$.
\item Let $M=\Spec k$ and give its line $L$ geometric fibre weight $-1$. Its nonzero section has weight $\Phi=1$, so $d=1$ is critical. The level twist cancels that fibre action, and the auxiliary origin again survives with stabiliser $\Gm$.
\end{enumerate}}

\subsection*{D--F}

\examplegloss{deligne-mumford-stack}{Deligne--Mumford stack}{
\begin{enumerate}[label=(\arabic*),leftmargin=*,itemsep=2pt,topsep=3pt]
\item The classifying stack $B\mu_2$ is DM because the map $\Spec k\to B\mu_2$ is an étale atlas. Its geometric object has two automorphisms, showing that DM does not mean absence of automorphisms.
\item The stack $[\A^1/\mu_2]$ is DM with étale atlas $\A^1$. It is a scheme away from the origin, but the origin retains $\mu_2$. Thus an étale atlas can coexist with stabilisers that vary from point to point.
\end{enumerate}}

\examplegloss{descent-datum}{Descent datum}{
\begin{enumerate}[label=(\arabic*),leftmargin=*,itemsep=2pt,topsep=3pt]
\item On $T=U\cup V$, take trivial lines on $U,V$ and identify them by the constant multiplier $1$ on $U\cap V$. Together with the inverse identification, this is descent data for the trivial line on $T$.
\item On the standard charts $U_0,U_1$ of $\PP^1$, the overlap coordinate $z=X_1/X_0$ is invertible. Multiplication by $z$ glues the two trivial lines into a line bundle. The transition isomorphism, rather than only the two local isomorphism classes, is essential data.
\end{enumerate}}

\examplegloss{descent-of-a-line-bundle-to-the-coarse-space}{Descent of a line bundle to the coarse space}{
\begin{enumerate}[label=(\arabic*),leftmargin=*,itemsep=2pt,topsep=3pt]
\item On $B\mu_2$, the trivial-character line is the pullback of the ordinary one-dimensional line on its coarse point. The stabiliser acts trivially, so descent is explicit.
\item The sign-character line on $B\mu_2$ does not descend: the coarse point has no automorphism that could induce the sign action. Its tensor square does descend, because the square of the sign character is trivial. The obstruction is the fibre representation.
\end{enumerate}}

\examplegloss{diagonal}{Diagonal}{
\begin{enumerate}[label=(\arabic*),leftmargin=*,itemsep=2pt,topsep=3pt]
\item For $\A^1$, the diagonal is the closed line $x=x'$ in $\A^2$. Pulling it back along two functions $f,g$ on $T$ gives their equality locus $f-g=0$.
\item Pull back the diagonal of $B\mu_2$ along its geometric object paired with itself. The resulting isomorphism space is $\mu_2$, with two points, because that object has two automorphisms. A stack diagonal therefore records more than equality of underlying coarse points.
\end{enumerate}}

\examplegloss{dimension-of-a-quotient-stack}{Dimension of a quotient stack}{
\begin{enumerate}[label=(\arabic*),leftmargin=*,itemsep=2pt,topsep=3pt]
\item The stack $[\A^1/\mu_2]$ has dimension $1-0=1$. A finite group has dimension zero, so passing to this finite quotient stack does not lower the dimension of the atlas.
\item The weighted line $\PP(2,1)$ has dimension $2-1=1$, because its atlas is $\A^2\setminus0$ and its presentation group is $\Gm$. Its finite boundary stabiliser does not change the subtraction by the dimension of that presentation group.
\end{enumerate}}

\examplegloss{effective-descent}{Effective descent}{
\begin{enumerate}[label=(\arabic*),leftmargin=*,itemsep=2pt,topsep=3pt]
\item Glue two trivial lines on an open cover by identity maps. The resulting global line is the trivial bundle; restricting it to the cover recovers exactly the original bundles and their identifications. This is effective descent.
\item The nontrivial torsor $s^2=t$ over $\mathbb G_{m,t}$ becomes trivial on its own étale cover. Its trivial local copies, equipped with the transition signs inherited from the torsor, glue back to that nontrivial torsor. Effectiveness allows the global object to be nontrivial even when every local object is trivial.
\end{enumerate}}

\examplegloss{equivariant-map}{Equivariant map}{
\begin{enumerate}[label=(\arabic*),leftmargin=*,itemsep=2pt,topsep=3pt]
\item Give the source $\A^1_z$ weight $1$ and the target $\A^1_w$ weight $2$. The map $w=z^2$ is equivariant, since $(tz)^2=t^2z^2$: applying the group before or after the map agrees.
\item The constant map from any $G$-scheme to a fixed point is equivariant. For example, $\A^1\to\{0\}\subset\A^1$ is equivariant for scaling. The same constant map to the nonfixed point $1$ is not equivariant.
\end{enumerate}}

\examplegloss{equivariant-vector-bundle}{Equivariant vector bundle}{
\begin{enumerate}[label=(\arabic*),leftmargin=*,itemsep=2pt,topsep=3pt]
\item Over the sign action on $\A^1_z$, give the trivial line the action $(z,v)\mapsto(-z,v)$. Its fibre at the origin is the trivial representation, and the corresponding stack line descends to the coarse line.
\item On the same underlying trivial line, use $(z,v)\mapsto(-z,-v)$. Now the origin fibre is the sign representation and descent fails. The underlying nonequivariant bundle has not changed; the equivariant structure has.
\end{enumerate}}

\examplegloss{etale-cover}{Étale cover}{
\begin{enumerate}[label=(\arabic*),leftmargin=*,itemsep=2pt,topsep=3pt]
\item The two standard opens $\A^1\sqcup\A^1\to\PP^1$ form an étale cover: open immersions are étale, and the two images cover the projective line.
\item The map $\mathbb G_{m,s}\to\mathbb G_{m,t}$, $t=s^2$, is a single finite étale cover. The derivative $2s$ is invertible, and every geometric $t\ne0$ has a square root. This cover supplies roots that cannot be chosen on any nonempty Zariski open of $\mathbb G_{m,t}$.
\end{enumerate}}

\examplegloss{etale-morphism}{Étale morphism}{
\begin{enumerate}[label=(\arabic*),leftmargin=*,itemsep=2pt,topsep=3pt]
\item The open immersion $\A^1\setminus\{0\}\hookrightarrow\A^1$ is étale. It introduces no equations or branching locally, although it is not surjective and therefore is not by itself a cover of $\A^1$.
\item The squaring map on $\Gm$ is étale because its derivative is invertible there. The same formula on $\A^1$ is not étale at zero: the derivative $2s$ vanishes. Restricting away from the branch point changes the answer.
\end{enumerate}}

\examplegloss{family}{Family}{
\begin{enumerate}[label=(\arabic*),leftmargin=*,itemsep=2pt,topsep=3pt]
\item Take $T=\A^1_a$ and the map $T\to\A^1_x$ defined by $x=a$. This is a family of affine-line points; its fibre at $a=c$ is the point $x=c$. A map $T'\to T$ changes the parameter by substitution.
\item The line bundle $\OO_{\PP^1}(1)$ is a family of one-dimensional vector spaces over $\PP^1$, hence a map $\PP^1\to B\Gm$. Every geometric fibre is a line, but the family is globally nontrivial: its transition $z=X_1/X_0$ cannot be removed by changing the two local frames. The units on either affine chart are constants.
\end{enumerate}}

\examplegloss{fibre-product-of-stacks}{Fibre product of stacks}{
\begin{enumerate}[label=(\arabic*),leftmargin=*,itemsep=2pt,topsep=3pt]
\item For ordinary schemes, $\A^1\times_{\Spec k}\Spec k=\A^1$.
There are no extra isomorphism choices, because scheme-valued objects
have only identity arrows.
\item For the atlas $\Spec k\to B\mu_2$, the fibre product
$\Spec k\times_{B\mu_2}\Spec k$ is $\mu_2$, consisting of two points.
The two points are the two isomorphisms between the trivial torsor and itself.
\end{enumerate}}

\examplegloss{flat-and-faithfully-flat}{Flat and faithfully flat}{
\begin{enumerate}[label=(\arabic*),leftmargin=*,itemsep=2pt,topsep=3pt]
\item The map $\A^1_z\to\A^1_u$, $u=z^2$, is faithfully flat.
Indeed, $k[z]$ is a free $k[u]$-module with basis $1,z$, and the
map is surjective, including over $u=0$.
\item The open immersion $\Gm\hookrightarrow\A^1$ is flat because
its ring map is localisation $k[z]\to k[z,z^{-1}]$. It is not
faithfully flat: its image omits the origin.
\end{enumerate}}

\examplegloss{fppf-topology}{Fppf topology}{
\begin{enumerate}[label=(\arabic*),leftmargin=*,itemsep=2pt,topsep=3pt]
\item The single map $\A^1_z\to\A^1_u$, $u=z^2$, is an fppf cover:
it is finite, flat, and surjective. It is not an étale cover, since
it is ramified at zero.
\item Restricting to $\Gm$ gives $s\mapsto s^2$, which is both
étale and fppf. Over this cover the unit $t$ on the base has the
chosen square root $s$.
\end{enumerate}}

\examplegloss{free-action}{Free action}{
\begin{enumerate}[label=(\arabic*),leftmargin=*,itemsep=2pt,topsep=3pt]
\item Multiplication of $\Gm$ on itself is free: $tx=x$ with
$x\ne0$ forces $t=1$. The quotient stack $[\Gm/\Gm]$ is
the ordinary one-point scheme.
\item The sign action of $\mu_2$ on $\Gm$ is free, since
$-z=z$ is impossible for $z\ne0$ in characteristic zero.
The same action on $\A^1$ is not free: zero has stabiliser $\mu_2$.
\end{enumerate}}

\examplegloss{functor-of-points}{Functor of points}{
\begin{enumerate}[label=(\arabic*),leftmargin=*,itemsep=2pt,topsep=3pt]
\item A morphism $T\to\A^1$ is exactly a global regular function
on $T$. Thus $\A^1(T)=\Gamma(T,\OO_T)$; pulling back a point
means pulling back that function.
\item For $T=\Spec k[\varepsilon]/(\varepsilon^2)$, the functions
$0$ and $\varepsilon$ give different $T$-points of $\A^1$.
Both restrict to the origin over $\Spec k$, so ordinary geometric
points alone do not distinguish them.
\end{enumerate}}

\subsection*{G--L}

\examplegloss{geometric-fibre-weight}{Geometric fibre weight}{
\begin{enumerate}[label=(\arabic*),leftmargin=*,itemsep=2pt,topsep=3pt]
\item On a line over a fixed point, the action $t\cdot v=t^2v$
has geometric fibre weight $2$. In the notes' section convention,
a nonzero constant section of this line has weight $-2$.
\item The sign line on $B\mu_2$ has action $v\mapsto-v$.
Its fibre character is $1$ modulo $2$; its tensor square has
trivial character, although the original line does not.
\end{enumerate}}

\examplegloss{geometric-point}{Geometric point}{
\begin{enumerate}[label=(\arabic*),leftmargin=*,itemsep=2pt,topsep=3pt]
\item The assignment $z\mapsto3$ defines a geometric point
$\Spec k\to\A^1_z$. Its only automorphism as a point of
the scheme is the identity.
\item The origin defines a geometric object of $[\A^1/\mu_2]$
for the sign action. It has two automorphisms, $1$ and $-1$.
Its image in the coarse line is an ordinary point without these arrows.
\end{enumerate}}

\examplegloss{git-quotient}{GIT quotient}{
\begin{enumerate}[label=(\arabic*),leftmargin=*,itemsep=2pt,topsep=3pt]
\item For the weight-$(1,1)$ action on $\A^2$, with geometric
fibre character $+1$ on the trivial linearised line, the semistable locus is $\A^2\setminus0$ and
the quotient is $\Proj k[x,y]=\PP^1$. Here both variables
have tensor degree one.
\item For weights $(2,1)$ the quotient stack is $\PP(2,1)$,
while the GIT quotient is $\Proj k[A,B]$ with degrees $(2,1)$.
Its second Veronese is $k[A,B^2]$, giving coarse space $\PP^1$.
\end{enumerate}}

\examplegloss{good-moduli-space}{Good moduli space}{
\begin{enumerate}[label=(\arabic*),leftmargin=*,itemsep=2pt,topsep=3pt]
\item The coarse map $[\A^1/\mu_2]\to\A^1_u$, $u=z^2$,
is also a good moduli space. Functions downstairs are the even
polynomials, and taking invariants is exact by averaging.
\item For scaling, $[\A^1/\Gm]\to\Spec k$ is good: functions
of weight zero are constants, and taking weight-zero parts is exact.
It identifies zero and the nonzero orbit, so it is not coarse.
\end{enumerate}}

\examplegloss{groupoid}{Groupoid}{
\begin{enumerate}[label=(\arabic*),leftmargin=*,itemsep=2pt,topsep=3pt]
\item A set $\{a,b\}$ becomes a groupoid by giving each object
only its identity arrow. It remembers two objects and no nontrivial
symmetries; this is a discrete groupoid.
\item A one-object groupoid with automorphism group $\mu_2$
has two arrows, identity and sign, with sign composed with itself
equal to identity. It has just one isomorphism class but retains a symmetry.
\end{enumerate}}

\examplegloss{groupoid-presentation}{Groupoid presentation}{
\begin{enumerate}[label=(\arabic*),leftmargin=*,itemsep=2pt,topsep=3pt]
\item Using the identity atlas of $\A^1$, the arrow space is
$\A^1\times_{\A^1}\A^1=\A^1$ and both projections are
identity. There is only one arrow at each object.
\item For $B\mu_2$, take object scheme $\Spec k$ and arrow
scheme $\mu_2$. Both source and target go to the one object,
and arrow composition is multiplication in $\mu_2$.
\end{enumerate}}

\examplegloss{inertia-stack}{Inertia stack}{
\begin{enumerate}[label=(\arabic*),leftmargin=*,itemsep=2pt,topsep=3pt]
\item For the scheme $\A^1$, every object has only the identity
automorphism. Consequently $I\A^1\simeq\A^1$: choosing an
automorphism adds no data.
\item For the sign quotient, the identity automorphism is available
everywhere, while sign fixes only zero. Hence
$I[\A^1/\mu_2]\simeq[\A^1/\mu_2]\sqcup B\mu_2$.
The second component records the extra automorphism at the origin.
\end{enumerate}}

\examplegloss{invariant-function-or-section}{Invariant function or section}{
\begin{enumerate}[label=(\arabic*),leftmargin=*,itemsep=2pt,topsep=3pt]
\item Under $z\mapsto-z$, the function $z^2$ is invariant:
$(-z)^2=z^2$. The function $z$ is not invariant, since it
changes sign.
\item Give the trivial line over $\A^1$ sign action on its fibre.
Then $f(z)=z$ is an equivariant section because $f(-z)=-f(z)$.
It defines a section of the line on the quotient stack despite
its coefficient function being noninvariant.
\end{enumerate}}

\examplegloss{isomorphism-sheaf}{Isomorphism sheaf}{
\begin{enumerate}[label=(\arabic*),leftmargin=*,itemsep=2pt,topsep=3pt]
\item Automorphisms of the trivial line on $T$ are multiplication
by invertible functions. Its isomorphism sheaf with itself is
$\Gm\times T$, rather than the additive line $\A^1\times T$:
zero multiplication is not invertible.
\item The trivial $\mu_2$-torsor over $T$ has automorphism sheaf
$\mu_2\times T$. It is the pullback of the diagonal of
$B\mu_2$ along the corresponding pair of identical families.
\end{enumerate}}

\examplegloss{line-bundle}{Line bundle}{
\begin{enumerate}[label=(\arabic*),leftmargin=*,itemsep=2pt,topsep=3pt]
\item The trivial line $\OO_{\A^1}$ has sections $k[z]$ and
one-dimensional fibres. Its nowhere-zero section $1$ gives a
global trivialisation.
\item The sign character defines a line on $B\mu_2$. Its pullback
to the atlas $\Spec k$ is an ordinary one-dimensional vector space,
but the retained sign action makes it a nontrivial stack line.
Its tensor square is trivial.
\end{enumerate}}

\examplegloss{linearisation}{Linearisation}{
\begin{enumerate}[label=(\arabic*),leftmargin=*,itemsep=2pt,topsep=3pt]
\item On $\A^1$ with $t\cdot z=tz$, lift the action to a trivial
line by $(z,v)\mapsto(tz,v)$. An invariant polynomial section
satisfies $f(tz)=f(z)$, so it is constant.
\item On the same underlying line use $(z,v)\mapsto(tz,tv)$.
Invariance now requires $f(tz)=t f(z)$, giving the multiples
of $z$. The bundle is unchanged, but its linearisation changes the sections.
\end{enumerate}}

\examplegloss{linearly-reductive-group}{Linearly reductive group}{
\begin{enumerate}[label=(\arabic*),leftmargin=*,itemsep=2pt,topsep=3pt]
\item For a $\mu_2$-representation with sign operator $\sigma$,
the map $v\mapsto(v+\sigma v)/2$ projects onto the invariant
vectors. Applying this averaging to a lift of an invariant vector
proves exactness of invariants.
\item A rational $\Gm$-representation splits into weight spaces
$V=\bigoplus_jV_j$. Its invariant space is the direct summand
$V_0$; an equivariant surjection is surjective on this summand.
\end{enumerate}}

\examplegloss{locally-free-sheaf}{Locally free sheaf}{
\begin{enumerate}[label=(\arabic*),leftmargin=*,itemsep=2pt,topsep=3pt]
\item The sheaf $\OO_{\A^1}^{\oplus2}$ is locally free of rank
two; the two constant basis sections trivialise it everywhere.
Its fibre at every $k$-point is $k^2$.
\item The line $\OO_{\PP^1}(1)$ is locally free of rank one:
each standard affine chart has a local generator, and the generators
are related by an invertible transition function. A global trivialisation
is not required for local freeness.
\end{enumerate}}

\subsection*{N--R}

\examplegloss{normal-line}{Normal line}{
\begin{enumerate}[label=(\arabic*),leftmargin=*,itemsep=2pt,topsep=3pt]
\item For the origin in $\A^1_z$, the conormal is $(z)/(z^2)$
and its dual is the normal line. It is an ordinary one-dimensional
vector space over $k$.
\item For $B\mu_2\subset[\A^1_z/\mu_2]$, the normal vector
changes sign. This normal line has sign character, whereas the
pullback of the coarse origin's ordinary normal line has trivial character.
\end{enumerate}}

\examplegloss{orbit}{Orbit}{
\begin{enumerate}[label=(\arabic*),leftmargin=*,itemsep=2pt,topsep=3pt]
\item Under sign on $\A^1$, the orbit of $1$ is $\{1,-1\}$,
while the orbit of zero is $\{0\}$. The first has trivial
stabiliser and the second has stabiliser $\mu_2$.
\item Under scaling by $\Gm$, all nonzero points form one orbit.
Its closure is the whole affine line and contains the separate
zero orbit. The invariant affine quotient identifies the two orbits.
\end{enumerate}}

\examplegloss{polarisation}{Polarisation}{
\begin{enumerate}[label=(\arabic*),leftmargin=*,itemsep=2pt,topsep=3pt]
\item On $\PP^1$, both $\OO(1)$ and $\OO(2)$ are ample
polarisations. Their spaces of sections have dimensions two and
three, respectively, although the underlying curve is the same.
\item On the lower-cut stack $\PP(2,1)$, the level-one polarisation
is $\OO(1)$ with sign character at the boundary. It does not
descend; its square $\OO(2)$ is the pullback of coarse $\OO_{\PP^1}(1)$.
\end{enumerate}}

\examplegloss{projective-degree}{Projective degree}{
\begin{enumerate}[label=(\arabic*),leftmargin=*,itemsep=2pt,topsep=3pt]
\item For $M=\A^1_x$ and $L=\OO_M$, write the section ring
as $k[x,h]$, where $\deg_m x=0$ and $\deg_m h=1$.
The function $x$ is a degree-zero element even though it is nonconstant.
\item If $x$ has group weight two and $h$ has weight zero,
then $xh$ has tensor degree one and group weight two;
$h^2$ has tensor degree two and group weight zero.
\end{enumerate}}

\examplegloss{pullback}{Pullback}{
\begin{enumerate}[label=(\arabic*),leftmargin=*,itemsep=2pt,topsep=3pt]
\item For any map $f:T\to\A^1$, the pullback
$f^*\OO_{\A^1}$ is $\OO_T$. The constant generator $1$
pulls back to the constant generator $1$.
\item For the coarse map $\pi:\PP(2,1)\to\PP^1$,
$\pi^*\OO_{\PP^1}(1)=\OO_{\PP(2,1)}(2)$.
Its boundary stabiliser acts trivially, because the fibre character
is $2=0$ modulo $2$.
\end{enumerate}}

\examplegloss{pushforward}{Pushforward}{
\begin{enumerate}[label=(\arabic*),leftmargin=*,itemsep=2pt,topsep=3pt]
\item For $\pi:[\A^1_z/\mu_2]\to\A^1_u$, $u=z^2$,
the structure-sheaf pushforward corresponds to the invariant module
$k[z]^{\mu_2}=k[u]$.
\item The sign line has pushforward module $z k[u]$, the odd
polynomials. This module is free of rank one over $k[u]$,
but the original stack line does not descend: its origin fibre
still has sign character.
\end{enumerate}}

\examplegloss{quasi-coherent-sheaf}{Quasi-coherent sheaf}{
\begin{enumerate}[label=(\arabic*),leftmargin=*,itemsep=2pt,topsep=3pt]
\item The module $k[z]$ gives the quasi-coherent sheaf
$\OO_{\A^1}$. On the open set $z\ne0$, its module of sections
is the localisation $k[z,z^{-1}]$.
\item The module $k[z]/(z)$ gives a quasi-coherent sheaf supported
at zero. Its fibre has dimension one at zero and zero away
from zero, so it is not a locally free sheaf on $\A^1$.
\end{enumerate}}

\examplegloss{quotient-stack}{Quotient stack}{
\begin{enumerate}[label=(\arabic*),leftmargin=*,itemsep=2pt,topsep=3pt]
\item For multiplication on $\Gm$, the quotient
$[\Gm/\Gm]\simeq\Spec k$. The torsor description recovers
an ordinary quotient because the action is free and transitive.
\item For the trivial action of $\mu_2$ on a point, the quotient
is $B\mu_2$. Over a test scheme it remembers a $\mu_2$-torsor;
over $k$ its one isomorphism class retains two automorphisms.
\end{enumerate}}

\examplegloss{regular-cutting-level}{Regular cutting level}{
\begin{enumerate}[label=(\arabic*),leftmargin=*,itemsep=2pt,topsep=3pt]
\item On $M=\A^1$ with point weight one and trivially
linearised $L=\OO_M$, the only fixed-component section weight
is zero. Thus $d=1$ is regular, and the lower cut is the
ordinary projective line.
\item Replace the point weight by two. The same $d=1$ is
still regular, but the lower cut is $\PP(2,1)$ with boundary
stabiliser $\mu_2$. Regularity allows this finite stabiliser.
\end{enumerate}}

\examplegloss{representable-morphism}{Representable morphism}{
\begin{enumerate}[label=(\arabic*),leftmargin=*,itemsep=2pt,topsep=3pt]
\item The atlas $\A^1\to[\A^1/\mu_2]$ is representable.
After a base change $T\to[\A^1/\mu_2]$, it becomes the
torsor $P\to T$, which is a finite étale morphism of schemes.
\item In contrast, $B\mu_2\to\Spec k$ is not representable:
base changing by $\Spec k$ leaves $B\mu_2$, whose nontrivial
automorphisms prevent it from being an algebraic space.
\end{enumerate}}

\examplegloss{residual-action}{Residual action}{
\begin{enumerate}[label=(\arabic*),leftmargin=*,itemsep=2pt,topsep=3pt]
\item On $\PP(2,1)$ the action $r\cdot[x:y]=[rx:y]$ commutes with the cutting action. On the interior coordinate $z=x/y^2$ it becomes $z\mapsto rz$. The point $z=0$ is fixed by this residual action, but has trivial stack stabiliser.
\item The boundary $[1:0]$ is also fixed as an isomorphism class under the residual action. Its stack stabiliser is $\mu_2$, coming from $t^2=1$ in the cutting group. These automorphisms and the residual symmetry are different data.
\end{enumerate}}

\examplegloss{residual-gerbe-at-a-closed-point}{Residual gerbe at a closed point}{
\begin{enumerate}[label=(\arabic*),leftmargin=*,itemsep=2pt,topsep=3pt]
\item At the origin of $[\A^1_z/\mu_2]$, the reduced closed substack is $B\mu_2$: it retains one point together with its identity and sign automorphisms. This is the residual gerbe of the origin.
\item At a nonzero coarse point of the same quotient, the scalar action is free. The residual gerbe is an ordinary $\Spec k$. The two cases differ through their automorphisms, although both coarse points are ordinary points of $\A^1$.
\end{enumerate}}

\examplegloss{root-divisor}{Root divisor}{
\begin{enumerate}[label=(\arabic*),leftmargin=*,itemsep=2pt,topsep=3pt]
\item On $\sqrt[2]{(\A^1_u,\{0\})}=[\A^1_z/\mu_2]$, the root divisor is cut out by $z=0$, so it is $B\mu_2$. Since $u=z^2$, the pullback of the ordinary origin divisor is twice this divisor.
\item On $\PP(2,1)$, the section $y$ vanishes simply along the root divisor $\cD=\{y=0\}$. The coarse boundary point has local parameter $y^2/x$, which vanishes to order two upstairs. Thus again $\pi^*D=2\cD$.
\end{enumerate}}

\examplegloss{root-gerbe}{Root gerbe}{
\begin{enumerate}[label=(\arabic*),leftmargin=*,itemsep=2pt,topsep=3pt]
\item Square roots of the trivial line on $\Spec k$ form $B\mu_2$. The root is a line with a specified square trivialisation. Multiplication by $-1$ preserves that trivialisation, giving a nonidentity automorphism.
\item Square roots of $\OO_{\A^1}$ form $\A^1\times B\mu_2$. Every geometric point, including nonzero points, has a $\mu_2$ stabiliser. This illustrates why omitting a root section leaves stabilisers everywhere, rather than only over a divisor.
\end{enumerate}}

\examplegloss{root-stack-of-a-divisor}{Root stack of a divisor}{
\begin{enumerate}[label=(\arabic*),leftmargin=*,itemsep=2pt,topsep=3pt]
\item Taking the square root of the origin divisor in $\A^1_u$ gives $[\A^1_z/\mu_2]$ with $u=z^2$. Away from zero the root section is nonzero and removes all automorphisms; at zero its sign ambiguity remains.
\item Taking the square root of one point $D\subset\PP^1$ gives $\PP(2,1)$. Its ordinary affine chart away from $D$ and cyclic quotient chart near $D$ show explicitly how a root stack modifies just the chosen divisor.
\end{enumerate}}

\subsection*{S--Z}

\examplegloss{scheme-theoretic-fibre}{Scheme-theoretic fibre}{
\begin{enumerate}[label=(\arabic*),leftmargin=*,itemsep=2pt,topsep=3pt]
\item For the scheme map $\A^1_z\to\A^1_u$, $u=z^2$, the fibre over zero is $\Spec k[z]/(z^2)$. It has one underlying point but retains the nonzero nilpotent class of $z$.
\item For the stack coarse map $[\A^1_z/\mu_2]\to\A^1_u$, base change gives $[\Spec k[z]/(z^2)/\mu_2]$. Its reduction is $B\mu_2$. Merely listing its geometric objects would miss the thickening.
\end{enumerate}}

\examplegloss{section-ring}{Section ring}{
\begin{enumerate}[label=(\arabic*),leftmargin=*,itemsep=2pt,topsep=3pt]
\item For $M=\PP^1$ and $L=\OO(1)$, sections of $L^m$ are degree-$m$ homogeneous polynomials in $X_0,X_1$. Thus the section ring is $k[X_0,X_1]$, with both variables in degree one.
\item For $M=\A^1_x$ and $L=\OO_M$, introduce its constant degree-one section $h$. Then $R=k[x,h]$, with $\deg x=0$ and $\deg h=1$. The nonconstant function $x$ belongs to the degree-zero algebra.
\end{enumerate}}

\examplegloss{section-weight}{Section weight}{
\begin{enumerate}[label=(\arabic*),leftmargin=*,itemsep=2pt,topsep=3pt]
\item If $t\cdot x=t^2x$ on $\A^1$ and $L=\OO$ has trivial fibre action, direct pullback gives $x(tp)=t^2x(p)$. The section $x$ has weight $2$, and $x^3$ has weight $6$.
\item Over a fixed point, give a line geometric fibre action $v\mapsto tv$. The section action includes the inverse fibre map, so its constant nonzero section has weight $-1$. This is the sign used for the level-one twist.
\end{enumerate}}

\examplegloss{semiprojective}{Semiprojective}{
\begin{enumerate}[label=(\arabic*),leftmargin=*,itemsep=2pt,topsep=3pt]
\item For $M=\A^1$, the affinisation is the identity $\A^1\to\Spec k[x]$. An isomorphism is projective, so $\A^1$ is semiprojective in the convention here, although it is not proper over $k$.
\item For $M=\PP^1$, all global regular functions are constant. Its affinisation is therefore the projective morphism $\PP^1\to\Spec k$. These two examples show how the convention includes both affine and projective varieties.
\end{enumerate}}

\examplegloss{semistable-point}{Semistable point}{
\begin{enumerate}[label=(\arabic*),leftmargin=*,itemsep=2pt,topsep=3pt]
\item In the lower cut with point weights $(2,1)$ and level $1$, $(0,1)$ is semistable because $y\one_1$ is nonzero there. The point $(1,0)$ is semistable because $x\one_1^2$ is nonzero there.
\item At level $0$, the constant twist section has positive tensor degree and weight zero, so it is nonzero even at $(0,0)$. That point is semistable although its stabiliser is the full group $\Gm$.
\end{enumerate}}

\examplegloss{separated-stack}{Separated stack}{
\begin{enumerate}[label=(\arabic*),leftmargin=*,itemsep=2pt,topsep=3pt]
\item The scheme $\A^1$, viewed as a stack, is separated: its diagonal in $\A^1\times\A^1$ is the closed subscheme $x=y$. A closed immersion is proper.
\item The stack $B\mu_2$ is also separated. After trivialising two torsors, its isomorphism sheaf is a copy of the finite group $\mu_2$. Its diagonal is finite and therefore proper, despite its nontrivial automorphisms.
\end{enumerate}}

\examplegloss{smooth-stack}{Smooth stack}{
\begin{enumerate}[label=(\arabic*),leftmargin=*,itemsep=2pt,topsep=3pt]
\item The stack $[\A^1/\mu_2]$ is smooth because $\A^1$ is a smooth étale atlas. Its coarse space is also a smooth affine line, with coordinate $z^2$.
\item The stack $[\A^2/\mu_2]$ for simultaneous sign change is smooth for the same atlas reason. Its coarse space is the singular surface $ac=b^2$, with $a=x^2$, $b=xy$, $c=y^2$. Smoothness of the stack does not force smoothness of its coarse space.
\end{enumerate}}

\examplegloss{stabiliser-or-isotropy-group}{Stabiliser or isotropy group}{
\begin{enumerate}[label=(\arabic*),leftmargin=*,itemsep=2pt,topsep=3pt]
\item At $z=0$ in $[\A^1_z/\mu_2]$, both $1$ and $-1$ fix the representative. The stabiliser is $\mu_2$, and it acts by sign on the tangent direction.
\item At any point of $\PP(2,1)$ with $y\ne0$, a stabilising $t$ must satisfy $ty=y$, hence $t=1$. Its stabiliser is trivial even though the group used in the quotient presentation is one-dimensional.
\end{enumerate}}

\examplegloss{stable-point}{Stable point}{
\begin{enumerate}[label=(\arabic*),leftmargin=*,itemsep=2pt,topsep=3pt]
\item For scalar multiplication, use the trivial line of geometric fibre character $+1$, whose semistable locus is $\A^2\setminus0$. Each orbit is a line with the origin removed. It is closed within this punctured plane and has trivial stabiliser. These points give the usual stable presentation of $\PP^1$.
\item For weights $(2,1)$ with the same fibre character $+1$, the orbit of $(1,0)$ is the nonzero $x$-axis. It is closed in $\A^2\setminus0$ and has stabiliser $\mu_2$. Thus this point is stable, although the action there is not free.
\end{enumerate}}

\examplegloss{stack}{Stack}{
\begin{enumerate}[label=(\arabic*),leftmargin=*,itemsep=2pt,topsep=3pt]
\item The scheme $\A^1$ gives a stack whose objects over $T$ are regular functions on $T$ and whose only arrows are identities. Compatible local functions glue, so ordinary schemes fit the stack axioms.
\item The stack $B\mu_2$ instead assigns to $T$ its groupoid of $\mu_2$-torsors. These torsors and their isomorphisms glue, while geometric objects retain sign automorphisms. The same gluing principle accommodates nontrivial symmetries.
\end{enumerate}}

\examplegloss{tame-stack}{Tame stack}{
\begin{enumerate}[label=(\arabic*),leftmargin=*,itemsep=2pt,topsep=3pt]
\item On $B\mu_2$, a vector space with involution $\sigma$ is pushed to its invariant part. The averaging operator $v\mapsto(v+\sigma v)/2$ makes taking invariants exact. This exhibits tameness over the coarse point.
\item On $[\A^1/\mu_2]$, the same averaging works for equivariant modules over $k[z]$. Their invariant modules form the coarse pushforward over $k[z^2]$, and exact sequences remain exact. Tameness does not remove the sign characters on stack line fibres.
\end{enumerate}}

\examplegloss{tensor-degree}{Tensor degree}{
\begin{enumerate}[label=(\arabic*),leftmargin=*,itemsep=2pt,topsep=3pt]
\item On $\PP^1$, the monomial $X_0^2$ is a section of $\OO(1)^{\otimes2}=\OO(2)$, so its tensor degree is $2$. This degree counts powers of the line bundle.
\item On the affine-line example with $t\cdot x=t^2x$, the element $x^7h^3$ lies in $H^0(\A^1,\OO^{\otimes3})$. Its tensor degree is $3$, while its group weight is $14$. These are independent indices.
\end{enumerate}}

\examplegloss{topology-in-descent-theory}{Topology, in descent theory}{
\begin{enumerate}[label=(\arabic*),leftmargin=*,itemsep=2pt,topsep=3pt]
\item The two standard affine charts of $\PP^1$ form a Zariski cover. A line bundle can be described by trivial lines on these charts and a transition function on their overlap.
\item The squaring map $\Gm\to\Gm$ is an étale cover. It allows a square root of the target coordinate after base change, although no such root exists on a nonempty Zariski open. Changing the topology changes which local choices are available for gluing.
\end{enumerate}}

\examplegloss{torsor}{Torsor}{
\begin{enumerate}[label=(\arabic*),leftmargin=*,itemsep=2pt,topsep=3pt]
\item The product $\Gm\times T\to T$, with multiplication on the first factor, is a trivial $\Gm$-torsor. Its identity section gives the chosen trivialisation; multiplying that section by a unit gives another one.
\item The map $\Spec k[s,s^{-1}]\to\Spec k[t,t^{-1}]$, $t=s^2$, is a $\mu_2$-torsor. It has no global section because $t$ has no square root in $k[t,t^{-1}]$, but it becomes trivial after pulling back along itself.
\end{enumerate}}

\examplegloss{torus}{Torus}{
\begin{enumerate}[label=(\arabic*),leftmargin=*,itemsep=2pt,topsep=3pt]
\item The group $\Gm$ is a one-dimensional torus. It acts on $\A^1$ by scalar multiplication, or by the weighted formula $x\mapsto t^2x$. Its characters are indexed by integers.
\item The product $\Gm\times\Gm$ is a two-dimensional torus, acting on $\A^2$ by $(t_1,t_2)\cdot(x,y)=(t_1x,t_2y)$. A character has the form $t_1^a t_2^b$, indexed by the pair $(a,b)\in\mathbb Z^2$.
\end{enumerate}}

\examplegloss{transporter}{Transporter}{
\begin{enumerate}[label=(\arabic*),leftmargin=*,itemsep=2pt,topsep=3pt]
\item For the sign action on $\A^1$, the transporter from $1$ to $-1$ is $\{-1\}$. It is the space of arrows between these two representatives of the same quotient object.
\item The transporter from $0$ to itself is all of $\mu_2$, because both group elements fix zero. Thus the same transporter construction gives isomorphisms between different representatives and automorphisms when the representatives coincide.
\end{enumerate}}

\examplegloss{unramified-morphism}{Unramified morphism}{
\begin{enumerate}[label=(\arabic*),leftmargin=*,itemsep=2pt,topsep=3pt]
\item The map $\mathbb G_{m,s}\to\mathbb G_{m,t}$, $t=s^2$, is unramified: the relative differential relation $2s\,ds=0$ forces $ds=0$ because $2s$ is invertible. This map is also flat, hence étale.
\item The closed inclusion $\{0\}\hookrightarrow\A^1$ is unramified, since a point has no relative differentials over $\A^1$. It is not flat and is not étale. Thus unramified alone does not mean that a morphism is an étale chart.
\end{enumerate}}

\examplegloss{vector-bundle}{Vector bundle}{
\begin{enumerate}[label=(\arabic*),leftmargin=*,itemsep=2pt,topsep=3pt]
\item The sheaf $\OO_{\A^1}^{\oplus2}$ is the trivial rank-two vector bundle. Its fibre at every $k$-point is $k^2$, and its two standard basis sections give a global frame.
\item On $B\mu_2$, take the representation $\sigma(v_1,v_2)=(v_1,-v_2)$. It gives a rank-two vector bundle with one trivial and one sign summand. Its nontrivial fibre action prevents descent to a vector bundle on the coarse point.
\end{enumerate}}

\examplegloss{veronese-subring}{Veronese subring}{
\begin{enumerate}[label=(\arabic*),leftmargin=*,itemsep=2pt,topsep=3pt]
\item For $k[u,v]$ with both variables in degree one, the second Veronese is generated by $u^2,uv,v^2$. After regrading these have degree one and satisfy $(uv)^2=u^2v^2$. Its projective spectrum is still $\PP^1$.
\item For $k[A,B]$ with degrees $(2,1)$, every even-degree monomial is a product of $A$ and $B^2$. The second Veronese is therefore $k[A,B^2]$, with both generators in regraded degree one, giving the coarse projective line of the toy cut.
\end{enumerate}}

\examplegloss{weight}{Weight}{
\begin{enumerate}[label=(\arabic*),leftmargin=*,itemsep=2pt,topsep=3pt]
\item The point action $x\mapsto t^2x$ gives tangent weight $2$ at the origin. With the direct-pullback action used here, the coordinate function $x$ also has weight $2$.
\item Give a line over a fixed point the geometric fibre action $v\mapsto t^3v$. Its fibre weight is $3$, but its constant section has weight $-3$ under the declared section convention. Naming the representation is essential when stating a weight.
\end{enumerate}}

\examplegloss{weighted-projective-stack}{Weighted projective stack}{
\begin{enumerate}[label=(\arabic*),leftmargin=*,itemsep=2pt,topsep=3pt]
\item With weights $(1,1)$ on $\A^2\setminus0$, the action is free and the quotient stack $\PP(1,1)$ is the ordinary projective line. Either nonzero coordinate can be normalised without a finite ambiguity.
\item With weights $(2,1)$, the quotient is $\PP(2,1)$. The chart $y\ne0$ is $\A^1$, while the boundary has stabiliser $\mu_2$. Its coarse space is again $\PP^1$, so the weights change stack geometry without changing this coarse curve.
\end{enumerate}}

\examplegloss{weight-space}{Weight space}{
\begin{enumerate}[label=(\arabic*),leftmargin=*,itemsep=2pt,topsep=3pt]
\item For $k[x]$ with $x\mapsto t^2x$ under direct pullback, the weight-$2a$ space is $k x^a$ for $a\ge0$. Every odd-weight space is zero. The invariant part is the weight-zero space $k$.
\item For the $\mu_2$-representation $k^2$ with action $(v_1,v_2)\mapsto(v_1,-v_2)$, the two coordinate axes are the trivial and sign weight spaces. Here the weight labels are residues modulo $2$.
\end{enumerate}}

\examplegloss{zariski-versus-etale-local}{Zariski versus étale local}{
\begin{enumerate}[label=(\arabic*),leftmargin=*,itemsep=2pt,topsep=3pt]
\item On $\mathbb G_{m,t}$, an equation $s^2=t$ has no solution on any nonempty Zariski open: it has no solution even in $k(t)$. Pulling back by the étale squaring cover supplies exactly such a root.
\item On the chart $x\ne0$ of $\PP(2,1)$, setting $x=1$ requires choosing $t$ with $t^2=x^{-1}$, which is possible étale locally and leaves a sign ambiguity. On $y\ne0$, setting $y=1$ instead uses the unique regular choice $t=y^{-1}$.
\end{enumerate}}

\end{document}
